\documentclass[10pt,a4paper]{article}
\usepackage{fontspec,xeCJK}
\setmainfont{Times New Roman}
\setCJKmainfont[AutoFakeBold=2,AutoFakeSlant=.2]{SimSun}
\setCJKmonofont{SimSun}
\usepackage[margin=20mm,top=22mm,bottom=22mm]{geometry}
\usepackage{paracol,graphicx,hyperref,xurl,fancyhdr,needspace,xcolor}
\hypersetup{unicode=true,hidelinks}\urlstyle{same}
\setlength{\columnsep}{8mm}\setlength{\parindent}{0pt}\setlength{\parskip}{5pt}
\setlength{\emergencystretch}{3em}\tolerance=2000
\widowpenalty=5000\clubpenalty=5000
\pagestyle{fancy}\fancyhf{}\fancyhead[L]{\small Benton · Piracy as Political Prism}\fancyfoot[C]{\thepage}\setlength{\headheight}{14pt}
\begin{document}

{\Large\bfseries Piracy as Political Prism: Sea Raiders and the Rise of the International Atlantic\par}\medskip

Lauren Benton. Department of History, Yale University, New Haven, CT, USA. Email: lauren.benton@yale.edu. Comparative Studies in Society and History (2026), 1–28. Published online: 28 September 2026.\par
Reading edition based on publisher HTML. Original PDF: 28 pages. All 124 original notes retained after their referring paragraphs. Page labels indicate paragraph starting pages. Original authors © 2026. CC BY 4.0.\par
\url{https://doi.org/10.1017/S0010417526100607}\par
\url{https://creativecommons.org/licenses/by/4.0/}\par
Keywords: piracy; Artigas; Atlantic history; Latin America; privateering; sovereignty; international law; constitutions; recognition; terrorism\par\bigskip
\Needspace{10\baselineskip}\section*{Abstract}
{\scriptsize\color{gray} [p. 1]}\ This article places maritime violence at the center of the history of international ordering. It views the transformation of the nineteenth-century Atlantic from the deck of one ship, a Baltimore-built privateer called the \textit{Irresistible.} The ship’s captain seized Spanish and Portuguese shipping under a privateering commission from the Oriental Republic. Sponsorship of privateers by the Oriental Republic and other new Latin American governments reflected an open-ended Atlantic constitutionalism that recognized sovereignty at multiple scales. That constitutional context formed the background to three cases involving the \textit{Irresistible} and its crew that reached the U.S. Supreme Court. The cases show how privateering politics opened to wide-ranging debates about the definition of states, war, and piracy, while preserving the indeterminacies of sovereignty in the new Atlantic system. That system, closing as it formed, was composed of new states pressured to align with powerful empires. Privateering disputes featured sharpening rhetoric about piracy as a crime under the law of nations while at the same time preserving a regime that protected low-level maritime violence. Situating the politics of piracy at the nexus of domestic and international orders, the article highlights the political uses and legal limits of the discourse of pirates as enemies of all and connects this history to an understanding of terrorism and states in the international order.\par
{\scriptsize\color{gray} [p. 1]}\ For much of the last two decades, historians of piracy have been caught in a cycle of myth-making and myth-busting.\nobreak\textsuperscript{1} One group of scholars, joined by popular writers eager to profit from the robust market for pirate lore, depicts maritime raiders as colorful rogues. They portray most pirates as inherently anti-authoritarian actors who possessed a penchant for experimenting with alternative forms of political solidarity. Histories in this vein highlight early modern pirates’ pacts to share booty and their rituals for overturning rule by despotic captains. Some accounts credit pirates with founding radically innovative republics or depict them as footloose maritime raiders who contributed to Enlightenment thought, aided democratic revolutions, and fueled anti-imperial movements. Campaigns against piracy become, in this telling, synonymous with other forms of state oppression.\nobreak\textsuperscript{2}\par{\small 1. My thanks to Jatin Dua and Andrew Shryock for the invitation to present the inaugural Sylva Thrupp lecture and to the forum discussants, Jatin Dua, Nathan Perl-Rosenthal, and Judith Scheele. I also wish to thank Arlyck, Kevin, Brett, Annabel, de Graaf, Beatrice, Douglas, Lawrence, Frega, Ana, Golove, David, Hulsebosch, Daniel, Mehta, Jay, Sarah Seo, and Frank Stewart for helpful comments. The title adapts the phrase “rise of the international” from Devetak, Richard and Dunn, Tim, eds., The Rise of the International (Oxford: Oxford University Press, 2024).\par}{\small 2. Prominent examples are Linebaugh, Peter and Rediker, Marcus, The Many-Headed Hydra: Sailors, Slaves, Commoners, and the Hidden History of the Revolutionary Atlantic (Boston: Beacon Press, 2000); Rediker, Marcus, Villains of All Nations (Boston: Beacon Press, 2005); Graeber, David, Pirate Enlightenment, or the Real Libertalia (New York: Farrar, Straus, and Giroux, 2023).\par}
{\scriptsize\color{gray} [p. 2]}\ A second group of historians has critiqued representations of pirates as political innovators, revealing their connections to established, if fluid, networks of patronage and power. These scholars show that most early modern raiders claimed to be not pirates but privateers—mariners licensed by legitimate sovereigns to seize enemy ships. Studies track how raiders engaged actively in what I have called “legal posturing” to color their actions as lawful and improve their chances of securing booty as property. Rather than fomenting revolution, most raiders cultivated powerful sponsors and aspired to positions as propertied heads of household in settled communities. Campaigns against piracy had various political objectives and creatively deployed rhetoric about pirates as enemies of all, but polities mainly sought to contain piracy through its criminalization under domestic law.\nobreak\textsuperscript{3}\par{\small 3. On legal posturing, see Benton, Lauren, A Search for Sovereignty: Law and Geography in European Empires, 1400–1900 (Cambridge: Cambridge University Press, 201), chapter 3. On the complex motives of sea raiders, see Hanna, Mark, Pirate Nests and the Rise of the British Empire, 1570–1740 (Chapel Hill: University of North Carolina Press, 2017); and Prange, Sebastian, “A Trade of No Dishonor: Piracy, Commerce, and Community in the Western Indian Ocean, Twelfth to Sixteenth Century,” American Historical Review 116, 5 (Dec. 2011): 1269–93 DOI: 10.1086/ahr.116.5.1269.\par}
{\scriptsize\color{gray} [p. 2]}\ Having participated in developing the second perspective, I have had reason to despair over the resilience of the first. No assemblage of evidence seems capable of dislodging the view that pirates were anti-authoritarian rogues routinely regarded as enemies of all. Nor does it appear to matter that the doctrinal force of the language of universal enmity changed over time.\nobreak\textsuperscript{4} The finding that campaigns to contain raiding relied more on actions under domestic law than on the exercise of universal jurisdiction has also not slowed the production of romanticized pirate histories.\nobreak\textsuperscript{5} Paradoxically, chiseling away at the image of the outcast pirate has perhaps helped focus attention on the figure of the sea raider rather than encouraging analysis of the institutional context that structured maritime relations.\par{\small 4. Vergerio, Claire, War, States, and International Order: Alberico Gentili and the Foundational Myth of the Laws of War (New York: Cambridge University Press, 2022), 108–118 ; cf. Heller-Roazen, Daniel, The Enemy of All: Piracy and the Law of Nations (Princeton: Princeton University Press, 2009).\par}{\small 5. Benton, Lauren, “Toward a New Legal History of Piracy: Maritime Legalities and the Myth of Universal Jurisdiction,” International Journal of Maritime History 23, 1 (2011): 1–15 ; Coakley, J., Kwan, C. N., and Wilson, D., “Introduction: Piracy and Occasional State Power,” International Journal of Maritime History 32, 3 (2020): 656–65 DOI: 10.1177/0843871420944651.\par}
{\scriptsize\color{gray} [p. 2]}\ Several lines of inquiry disrupt these debates by underscoring the political significance of maritime raiding. One strand explores the relation of maritime violence to state formation. Historians have tracked the effects of piracy on political authority in such diverse contexts as the ancient Mediterranean, sixteenth-century South Asia, and the early nineteenth-century United States.\nobreak\textsuperscript{6} A closely related strand of research investigates the place of piracy in the rise of regional and global orders.\nobreak\textsuperscript{7} Lisa Ford and I have traced, for example, how imperial and interimperial maritime campaigns to contain piracy helped produce gap-ridden, jurisdictionally complex legal regimes of oceanic dimensions.\nobreak\textsuperscript{8}\par{\small 6. de Souza, Philip, “Rome’s Contribution to the Development of Piracy,” Memoirs of the American Academy in Rome. Supplementary Volumes, vol. 6, 2008, 71–96 ; Subramanian, Lakshmi, The Sovereign and the Pirate: Ordering Maritime Subjects in India’s Western Littoral (New Delhi, Oxford University Press, 2016) DOI: 10.1093/acprof:oso/9780199467044.001.0001; Wang, Wensheng, White Lotus Rebels and South China Pirates: Crisis and Reform in the Qing Empire (Cambridge: Harvard University Press, 2014); and Arlyck, Kevin, The Nation at Sea: The Federal Courts and American Sovereignty, 1789–1825 (Cambridge: Cambridge University Press, 2025) DOI: 10.1017/9781009393041.\par}{\small 7. This approach was pioneered by Janice E. Thomson in \textit{Mercenaries, Pirates, and Sovereigns} (Princeton: Princeton University Press, 1996). For an overview of this perspective, see Renaud Morieux and Jeppe Mulich, eds., “Ordering the Oceans, Ordering the World: Law, Violence, and European Empires,” \textit{Past \& Present} 265, Supplement 17 (November 2024): 1–20.\par}{\small 8. Benton, Lauren and Ford, Lisa, Rage for Order: The British Empire and the Origins of International Law (Harvard University Press, 2016) DOI: 10.4159/9780674972780, chapter 5; Lisa Ford and Lauren Benton, “Empire, Nation, and the International in the Mid-Century Atlantic,” in Devetak and Dunne, eds, \textit{The Rise of the International}, 189–209.\par}
{\scriptsize\color{gray} [p. 3]}\ This article builds on legal-political approaches to piracy to examine maritime raiding and the rise of the international Atlantic. Setting aside older debates about pirates centered on piracy’s so-called golden age, I analyze the connections between a surge in Atlantic raiding and the proliferation of states in the early nineteenth century. The inquiry requires some critical distance from narratives about the Age of Revolutions as a period in which a “wave” of revolutionary ideology or a “contagion of sovereignty” transformed the region.\nobreak\textsuperscript{9} As Christopher Clark notes in relation to the global political upheavals of 1848, revolutions did not always respond to global currents but instead developed as “cognate” events that unfolded “within kindred cultural and political orders.”\nobreak\textsuperscript{10} In the early nineteenth-century Atlantic, new political entities formed in sync in response to imperial tensions and war, in particular the crisis created by the French invasion of the Iberian Peninsula. An open-ended constitutionalism framed a variety of possible political formations, and new entities grappled with imperial governments that sought in some ways to foster and in other ways to obstruct or limit the formation of new sovereignties.\par{\small 9. The “contagion of sovereignty” is from Armitage, David, The Declaration of Independence: A Global History (Cambridge: Harvard University Press, 2007), 103–4 DOI: 10.4159/9780674020276-004. An example of overreliance on the metaphor of waves is Sivasundaram, Sujit, Waves Across the South: A New History of Revolution and Empire (Chicago: University of Chicago Press, 2021) DOI: 10.7208/chicago/9780226790558.001.0001.\par}{\small 10. Clark, Christopher, Revolutionary Spring: Europe Aflame and the Fight for a New World, 1848–1849 (New York: Crown, 2023), 279 .\par}
{\scriptsize\color{gray} [p. 3]}\ There was also more to the region’s political transformation than a set of cognate events. A perceived urgency to establish claims to international belonging gripped the region. Diplomats, merchants, officials, mariners, laborers, and captives surveyed international events and pressures with intense interest and positioned themselves in relation to these changes. Explaining political reordering on an oceanic scale in the period requires close attention to the interpenetration of domestic and international political arenas.\nobreak\textsuperscript{11} It depends, that is, on treating the international order as something other than an already constituted realm hovering beyond the horizon of local, national, and regional politics.\nobreak\textsuperscript{12}\par{\small 11. On the relation of inside/outside orders in political thought, see Brett, Annabel, Changes of State: Nature and the Limits of the City in Early Modern Natural Law (Princeton: Princeton University Press, 2011); on inside/outside dimensions of empires, see Benton and Ford, \textit{Rage for Order}, 193–195.\par}{\small 12. Bartelson, Jens, Becoming International (New York: Cambridge University Press, 2023) DOI: 10.1017/9781009400718, chapter 1.\par}
{\scriptsize\color{gray} [p. 4]}\ This perspective shifts attention away from questions about whether pirates were insiders or outsiders. In this period as in others, the key characteristic of maritime raiding was its flickering quality, the fundamental instability of its legal protections.\nobreak\textsuperscript{13} Like constitutionalism itself, legal practices of privateering and piracy opened simultaneously toward domestic (national or imperial) and international ordering.\nobreak\textsuperscript{14} Just as disputes inside polities altered the conditions under which booty could be turned into property, imperial politics informed judgments about which states might be recognized as legitimate sponsors of raiding. Multi-sited conflicts about sea raiding transformed an Atlantic of proliferating political experiments into a regional system of unequal states under imperial sway.\par{\small 13. On these instabilities in the twentieth century, see Jatin Dua, \textit{Captured at Sea: Piracy and Protection in the Indian Ocean} (Oakland: University of California Press, 2019).\par}{\small 14. On the international dimensions of constitutionalism, see Golove, David M. and Hulsebosch, Daniel J., “A Civilized Nation: The Early American Constitution, the Law of Nations, and the Pursuit of International Recognition,” Columbia Law Review, 85, 4 (2010): 932–1066 ; LaCroix, Alison, The Interbellum Constitution: Union, Commerce, and Slavery in the Age of Revolutions (New Haven: Yale University Press, 2024). On American sovereignty forming in relation to international order, see Arylck, \textit{The Nation at Sea}; Gould, Eliga H., Among the Powers of the Earth: The American Revolution and the Making of a New World Empire (Cambridge: Harvard University Press, 2014); and Hendrickson, David, Peace Pact: The Lost World of the American Founding (Lawerence: University of Kansas Press, 2007).\par}
{\scriptsize\color{gray} [p. 4]}\ Privateering disputes were at the heart of this process. To trace their effects, this article views the changes in the Atlantic order from the deck of one ship. It follows the \textit{Irresistible} and its enterprising Baltimore captain, John Daniels, on a pair of cruises to capture Iberian vessels between 1818 and 1822.\nobreak\textsuperscript{15} Daniels sailed with a privateering commission from the Oriental Republic, a confederation led by the Banda Oriental, also known as the Provincia Oriental, a former territory of the Spanish empire in the Río de la Plata region. After Daniels returned to Baltimore with his booty, Portuguese and Spanish consuls challenged the captures, and Daniels also faced criminal prosecution. Scores of crew members of the \textit{Irresistible} who had joined a mutiny and led raids were tried and convicted of piracy. Three cases stemming from the captures of the \textit{Irresistible}—two prize cases and one criminal case—reached the U.S. Supreme Court and stimulated debates about much larger issues than the question of whether Daniels could keep his booty: What was a state? What was the difference between war and rebellion? What was the legal definition of piracy?\par{\small 15. Because he was referred to in U.S. court papers involving the \textit{Irresistible} as “John Daniels,” I use that name. He appears in some U.S. sources as John Daniel Danels. In Portuguese, Spanish, and South American sources, variations of the name include Danelly, Danele, and Danell. Two articles describing Daniels’s activities use “Danels”: Hopkins, Fred, “For Flag and Profit: The Life of Commodore John Daniel Danels of Baltimore,” Maryland Historical Magazine, 80:4 (1985); and Head, David, “Independence on the Quarterdeck: Three Baltimore Seafarers, Spanish America, and the Lives of Captains in the Early American Republic,” The Northern Mariner XXIII (January 2013): 1–20 .\par}
{\scriptsize\color{gray} [p. 5]}\ These questions emerged at a pivotal time in the formation of the international Atlantic. Following the \textit{Irresistible}, we spy an international order closing as it formed. When Daniels began his cruise, the region was crowded with political experiments, including juntas, city-states, confederations, and self-proclaimed nations. When the U.S. Supreme Court issued rulings about the \textit{Irresistible} and its captures, only a handful of new polities had acquired the status of nation-states. International recognition had become an offer that political communities and their leaders could not refuse.\nobreak\textsuperscript{16} Yet even as they received a measure of recognition, new states struggled to define and expand their sovereign capacities, and imperial governments retained enormous influence over this process. The new Atlantic system was a system of sovereign states with formal equality but vast differences in the capacity and right to make war.\par{\small 16. Bartelson, \textit{Becoming International}, 10.\par}
{\scriptsize\color{gray} [p. 5]}\ The emerging Atlantic order comprised a regime of limited violence that preserved a space for sea raiding. Empires leaned on American states to contain privateering, but few advocated banning it completely. The policy of slow, partial recognition of insurgent governments meant that the British government had to refrain from classing all South American raiders as pirates. The Spanish government, meanwhile, exerted strong pressure on U.S. officials to curb privateering while also sponsoring corsair raids of its own, particularly in the waters around Cuba. The result of these trends was that rhetoric about piracy as a crime under the law of nations was becoming sharper at the same time that protections of low-level maritime violence were being carefully preserved and revised. Becoming international in the Atlantic meant achieving membership in a stratified interpolity order under interimperial regulation and subject to the continued threat of imperial violence.\nobreak\textsuperscript{17}\par{\small 17. On global regimes of violence, see Benton, Lauren, They Called It Peace: Worlds of Imperial Violence (Princeton: Princeton University Press, 2024), 8 .\par}
{\scriptsize\color{gray} [p. 5]}\ Atlantic sea raiding and its regulation also brought some unintended consequences. Slave traders obtained and used privateering licenses, and counterrevolutionaries borrowed arguments about insurgent states’ right to make war. Like the capacious use of “terrorism,” the label of piracy was deployed mainly to degrade the status of sponsors of sea raiding within the international order rather than to position groups outside it. Despite claims about its foundational legal significance, ringing rhetoric about pirates as enemies of all derived its meaning from the politics of international ordering.\par
\Needspace{10\baselineskip}\section*{Atlantic Constitutionalism}
{\scriptsize\color{gray} [p. 5]}\ Our story begins not in Baltimore, the home port of John Daniels, but on the banks of the Río de la Plata, where a fluid constitutionalism framed responses to the collapse of the Spanish monarchy after the French invasion of the Iberian Peninsula. When a junta formed in Buenos Aires in 1810 to fill the political void, the scope of its authority across former Spanish territories remained uncertain. A set of provincial leaders and their followers were reluctant to trade rule by European empires for the dominance of port elites, and all groups were responding to the persisting threat of imperial attack and annexation.\nobreak\textsuperscript{18}\par{\small 18. Chiaramonte, José Carlos, Ciudades, provincias, estados: Orígenes de la Nación Argentina (1800–1846) (Buenos Aires: Espasa Calpe, 1997); Rodriguez, Jaime E., The Independence of Spanish America (Cambridge: Cambridge University Press, 1998) DOI: 10.1017/CBO9780511721137; and Echeverri, Marcela and Soriano, Cristina, eds., The Cambridge Companion to Latin American Independence (Cambridge: Cambridge University Press, 2023) DOI: 10.1017/9781108679336.\par}
{\scriptsize\color{gray} [p. 6]}\ One result was political fragmentation. On land bordering the Río Uruguay and the Río de la Plata, a war-tested colonel named José Gervasio Artigas joined the fight against loyalist Spanish forces concentrated at Montevideo, then engineered a break with Buenos Aires. Even before the 1815 rupture, Artigas was laying the groundwork for the territory’s semi-autonomy. At the beginning of 1813, when the second triumvirate to rule Buenos Aires called a Constituent Assembly, Artigas issued a set of instructions to delegates from the Banda Oriental. The instructions provide a primer of Artiguista thought—a constitutionalism that left open the possibility of sovereignty at multiple scales.\par
{\scriptsize\color{gray} [p. 6]}\ Artigas assigned the government of the United Provinces a limited set of capacities and proposed reserving some sovereign authority for the Banda Oriental. The United Provinces would assume key prerogatives, including the right to establish “fundamental laws” regulating relations among provinces; the right to represent the provinces in most—but not all—dealings with foreign powers; and the right to organize collective defense and maintain provincial military forces. Artigas supported these claims by stating that sovereignty resided in the people (\textit{los pueblos}) and was delegated, in different measure, to his government and to the government of the United Provinces.\nobreak\textsuperscript{19}\par{\small 19. Frega, Ana, “La virtud y el poder: La soveranía particular de los pueblos en el proyecto artiguista,” in Goldman, Noemi and Salvatore, Ricardo, eds., Caudillismos rioplatenses: Nuevas miradas a un viejo problema (Buenos Aires: Eudeba, 1998), 101–133 .\par}
{\scriptsize\color{gray} [p. 6]}\ Artigas’s instructions to the delegates from the Provincia Oriental drew from eclectic sources. The twenty articles include text paraphrasing passages from the U.S. Declaration of Independence, the Articles of Confederation, the Constitutions of Virginia and Massachusetts, and the U.S. Constitution.\nobreak\textsuperscript{20} The document resembled a declaration of independence in that it announced the dissolution of ties between the colonies and the Spanish monarchy and established the existence of the Provincia Oriental as a sovereign entity.\nobreak\textsuperscript{21} Paraphrasing the U.S. Articles of Confederation and the Constitution of Massachusetts, Article 11 stated that the Provincia Oriental would “retain its sovereignty, liberty, and independence, all power, jurisdiction, and rights not delegated explicitly by the Confederation of the United Provinces at the Congress.”\nobreak\textsuperscript{22} Article 10 characterized the pact among provinces as “a firm bond of friendship.”\nobreak\textsuperscript{23}\par{\small 20. These texts appeared as appendices of Thomas Paine’s “The Independence of the Mainland Justified.” On Artigas’s use of these sources, see Muñoz, Eugenio Petit, Artigas y su ideario a través de seis series documentales (Montevideo: Universidad de la República Oriental del Uruguay, 1956). See also Miranda, Hector, Las Instrucciones del año XIII (Montevideo: A. Barreiro y Ramos, 1910); Caetano, Gerardo and Ribeiro, Ana, eds., Las Instrucciones del año XIII: 200 años después (Montevideo: Planeta, 2013).\par}{\small 21. On declarations of independence as a distinctive genre, see David Armitage, \textit{The Declaration of Independence.}\par}{\small 22. Miranda, \textit{Las Instrucciones del año XIII}, 413. My translation.\par}{\small 23. Miranda, \textit{Las Instrucciones del año XIII}, 411–417. My translation.\par}
{\scriptsize\color{gray} [p. 7]}\ Around the time that he issued the instructions, Artigas explained that a variety of constitutional designs would be consistent with the recognition that sovereignty resided in the people. Addressing his followers, Artigas proclaimed that his own authority “emanates from you and ceases with your sovereign presence.”\nobreak\textsuperscript{24} The authority of the Constituent Assembly, he stated, did not flow from the transfer of the people’s loyalty from the Spanish crown to the new government. The assembly’s legitimacy could derive only from a pact that left the liberty of the provinces intact, and the foundation of any new political regime would also depend on the authority of \textit{los pueblos}. Artigas styled himself the \textit{Protector de} \textit{los Pueblos Libres} and explained to one delegate, “The particular sovereignty of the \textit{pueblos} will be explicitly declared and sustained as the singular object of our revolution.”\nobreak\textsuperscript{25} Although in theory the sovereignty of free peoples might be assigned to authorities at any level, Artigas regarded towns as political entities with special legitimacy. The phrase “pueblos libres” referred unambiguously to the sovereignty of the people (\textit{los pueblos}), but it also evoked the view that the smallest settlements (towns or \textit{pueblos}) were repositories of sovereignty close to the people—an idea that represented, as one scholar noted decades ago, “the essence of Artiguista thought.”\nobreak\textsuperscript{26}\par{\small 24. Speech of 4 April 1813 to the Congreso de Abril. José Gervasio Artigas and Lucía Sala de Touron, \textit{José Gervasio Artigas: Obra Selecta} (Caracas: Fundación Biblioteca Ayacuch, 2000), 22. My translation.\par}{\small 25. Quoted in Juan Antonio Rebelle, \textit{Lo hispánico en la acción y en el pensamiento político de Artigas} (Montevideo: s.n., 1953), 16–17. My translation.\par}{\small 26. Rebelle, \textit{Lo hispánico en la acción y en el pensamiento político de Artigas}, 20 (“la esencia del pensamiento artiguista”). The “proyecto de constitución para la Provincia Oriental” of 1813 begins by listing the “pueblos … que forman la Provincia Oriental”: Canelones, Piedras, Pando, Minas, Maldonado, San Carlos, Rocha, Santa Teresa, Santa Lucía, Pintado, San José, el Colla, Colonia, Espinillo, Víboras, Santo Domingo, Soriano, Mercedes, Porongos, Paysandú, Cerro Large, and Belén. Felipe Pigna, ed., \textit{José Gervasio Artigas: Escritos politicos y económicos} (Buenos Aires: Grupo Editorial Planeta, 2010), 51–52. On the central role of towns and cities in the Río de la Plata revolution, see Chiaramonte, \textit{Ciudades, provincias, estados}, 145–47; and Adelman, Jeremy, Sovereignty and Revolution in the Iberian Atlantic (Princeton: Princeton University Press, 2009).\par}
{\scriptsize\color{gray} [p. 7]}\ Artigas was not founding a nation.\nobreak\textsuperscript{27} Nor did he intend, with his extensive use of U.S. constitutional documents, merely to adapt a U.S. model of statehood.\nobreak\textsuperscript{28} That model was, in any case, full of ambiguities and not easy to appropriate. The United States was a hotbed of constitutional conflicts, not yet organized around the dyad of federal and state governments but involving a jumble of claimants to authority, including “towns, cities, and assemblies of private citizens.”\nobreak\textsuperscript{29} The country was also entering a slow-burning constitutional crisis centered on the extension of slavery to western territories.\nobreak\textsuperscript{30} Artigas was drawn to U.S. texts precisely because they affirmed the possibility of multiple configurations of sovereignty, flexibility that was was necessary to his political survival in the face of imperial aggression.\par{\small 27. Frega points out that Artigas consolidated his authority over the territory through state-like institutions (Ana Frega, “La virtud y el poder”).\par}{\small 28. Portillo, \textit{Crisis Atlantica}, 116.\par}{\small 29. LaCroix, \textit{The Interbellum Constitution}, 2.\par}{\small 30. Belko, William S., ed., Contesting the Constitution: Congress Debates the Missouri Crisis, 1819–1821 (Columbia: University of Missouri Press, 2021); Pasley, Jeffrey L. and Hammond, John Craig, eds., A Fire Bell in the Past: The Missouri Crisis at 200 (Columbia: University of Missouri Press, 2021).\par}
{\scriptsize\color{gray} [p. 7]}\ Multiple threats were looming. Artiguistas faced the likelihood of invasion by Luso-Brazilian forces, the danger of re-conquest by a resurgent Spanish empire, and the possibility that the British, who had invaded in 1806–1807, might take advantage of regional strife to mount another attempt at occupation. These conditions heightened interest in achieving international recognition.\nobreak\textsuperscript{31} In service to that goal and to fortify the polity’s ability to defend itself, Artigas claimed the right to make war, including the authority to sponsor privateering.\par{\small 31. On efforts of Latin American agents in London to establish diplomatic relations, see Racine, Karen, Francisco de Miranda, a Transatlantic Life in the Age of Revolution (Wilmington: Scholarly Resources, 2003); and Jaksic, Ivan, Andrés Bello: Scholarship and Nation-Building in Nineteenth-Century Latin America (New York: Cambridge University Press, 2001).\par}
\Needspace{10\baselineskip}\section*{Privateering in Atlantic Politics}
{\scriptsize\color{gray} [p. 8]}\ Artigas looked out on an ocean of raiding. By the second half of the eighteenth century, conventions for adjudicating captures of ships in wartime were relatively consistent across the Atlantic world. Each empire had its own prize courts, and captors were supposed to sail captured ships and their cargoes to their own courts for adjudication. In theory, all claimants, regardless of nationality, could present evidence that ships or cargoes had been seized unlawfully and should be returned.\nobreak\textsuperscript{32}\par{\small 32. Nathan Perl-Rosenthal, “An Interpolity Legal Regime in the Eighteenth Century: Procedural Law of Prize,” \textit{Past \& Present} 265, Supplement 17 (November 2024): 169–192.\par}
{\scriptsize\color{gray} [p. 8]}\ This relatively stable interpolity framework coexisted with a jumble of evasive practices. Ship owners, merchants, and mariners maneuvered to sidestep legal requirements or conform to them only on paper. Sea raiders creatively interpreted the terms of their commissions, and they sometimes did not take their captures to prize courts, or they presented their booty in courts of uncertain jurisdiction where captures could be adjudicated quickly and without close scrutiny of evidence. Captains and their sponsors were equally adept at legal maneuvering. It was possible to purchase privateering commissions and to alter the nationality of ships and goods by changing owners on paper or by transferring goods to neutral vessels before entering ports. A surge in enemy ships sailing as neutrals confounded the system, with international coalitions favoring different legal approaches to neutral shipping in wartime.\nobreak\textsuperscript{33} Incentives to profit from prizes stimulated the proliferation of prize courts around the region, with some remaining active in periods of relative interimperial peace.\par{\small 33. Kulsrud, Carl, Maritime Neutrality to 1780: A History of the Main Principles Governing Neutrality and Belligerency to 1780 (Boston: Little, Brown, 1936); Moats, Sandra, Navigating Neutrality: Early American Governance in the Turbulent Atlantic (Charlottesville: University of Virginia Press, 2021) DOI: 10.52156/m.5645; and Arlyck, \textit{The Nation at Sea}, Part I.\par}
{\scriptsize\color{gray} [p. 8]}\ Recognizing opportunities in this fluid system, new Latin American governments quickly began to license privateers and dispatch recruiters to North American ports. In 1816, American citizen Thomas Taylor arrived in Baltimore as an envoy of the Buenos Aires government carrying a small number of blank commissions, which he either gave or—much more likely—sold to Baltimore mariners. Other South American agents followed, and news surged through South and North American ports about opportunities to raid for new governments against Spanish shipping.\nobreak\textsuperscript{34}\par{\small 34. Bemis, Samuel Flagg, “Early Diplomatic Missions from Buenos Aires to the United States, 1811–1824,” Proceedings of the American Antiquarian Society, 49, 1 (April 1939): 29–30 .\par}
{\scriptsize\color{gray} [p. 9]}\ In the Río de la Plata, privateers were vital to defense. In 1815, after his forces had expelled Buenos Aires troops from Montevideo, Artigas led the movement to form the Liga Federal (Federal League), a confederation of the Banda Oriental and the provinces of Córdoba, Corrientes, Entre Ríos, Misiones, and Santa Fe. In 1816, he began to issue \textit{patentes de corso} (privateering commissions) listing Artigas as “the Chief of the Orientales and Protector of the Peoples of the República Oriental” and instructing holders to “respect the Flag of the Republic and attack its enemies in the manner judged most opportune according to the Law of Nations.”\nobreak\textsuperscript{35} The use of the “República Oriental” on the privateering commissions was intentional and strategic. It identified the confederation as an international entity with the right to make war and announced Artigas as a head of state. The practice was consistent with Artigas’s use of “República Oriental” only in external communications and of “Banda Oriental” and “Provincia Oriental” on documents circulating within the territory and region.\nobreak\textsuperscript{36} The bifurcation could also work against him. Luso-Brazilian officials would insist on representing Artigas as the head of the Banda Oriental, not the confederation, and the more limited designation was adopted by U.S. courts and U.S. government officials.\par{\small 35. Beraza, \textit{Los Corsarios de Artigas}, 32. My translation.\par}{\small 36. Frega, Ana, “Uruguayos y orientales: Itinerario de una síntesis compleja,” in Chiaramonte, J. C., Marichal, C., \& Granados, A., eds., Crear la nación. Los nombres de los países de América Latina: Identidades políticas y nacionalismo (Buenos Aires: Editorial Sudamericana, 2008), 95–112 .\par}
{\scriptsize\color{gray} [p. 9]}\ Captains quickly learned about the peculiar value of the commissions signed by Artigas. Whereas commissions from the government of the United Provinces specified that privateers could attack only Spanish ships, Artigas’s \textit{patentes} authorized captures of both Portuguese and Spanish vessels, opening to privateers the possibility of tapping into the lucrative trade from Brazil and Asia to Portugal. But there were downsides and dangers. The commissions were not accepted in Buenos Aires prize courts, and it remained uncertain whether courts in other jurisdictions, including the United States, would recognize the \textit{patentes} as legitimate authorization to capture enemy ships and cargo.\nobreak\textsuperscript{37} The disruptions to Portuguese-flagged shipping hardened Portuguese resolve to mount an invasion, and Montevideo fell to Luso-Brazilian forces in January 1817, followed by Colonia. The Oriental Republic persisted as a privateering sponsor, but it became a sponsor with no major port and no functioning prize court.\nobreak\textsuperscript{38}\par{\small 37. Agustín Beraza, \textit{Los Corsarios de Artigas} (Montevideo: Imprenta Nacional, 1949), 36; Bealer, Lewis Winkler, Los corsarios de Buenos Aires, sus actividades en las guerras hispano-americanas de la independencia, 1815–1821 (Buenos Aires: Publicaciones del Instituto de Investigaciones Históricas LXXII, 1937).\par}{\small 38. On Artigas sponsoring privateers without a port, see John Quincy Adams, Secretary of State, to Luis de Onís, in William R. Manning, \textit{Diplomatic Correspondence of the United States Concerning Independence of the Latin-American Nations, Vol. 1} (New York, Oxford University Press, 1925), 88–91, 89.\par}
{\scriptsize\color{gray} [p. 9]}\ There is no question that Artigas was deploying the commissions to position himself as a head of state. He burnished the official nature of the act of issuing the commissions by producing a document called the “Reglamento del Corso,” which outlined procedures for privateers and the adjudication of prizes. In 1817, in an effort to increase privateering attacks on Portuguese shipping in the midst of the Luso-Brazilian invasion, Artigas supplied a set of blank commissions to the United States consul in Buenos Aires, Thomas Lloyd Halsey, who met with Artigas at his headquarters in Purificación. There, Artigas gave Halsey a letter for President James Monroe, composed as a message from one head of state to another.\nobreak\textsuperscript{39}\par{\small 39. General José Artigas to James Monroe, 1 September 1817, in Manning, \textit{Diplomatic Correspondence}, 2175. Officials in Buenos Aires were furious that Halsey had met with Artigas. An order to expel Halsey from the city in January 1818 was quickly rescinded, but Pueyrredón, Supreme Director of the United Provinces, wrote to President James Monroe to complain that Halsey’s behavior had “bordered on the abuse of power” when he “entered into contracts concerning privateering which must necessarily compromise the Provinces under my command with the other Nations.” Juan Martín Pueyrredón, Supreme Director of the Untied Provinces of South America, to James Monroe, President of the United States, 31 January 1818, Manning, \textit{Diplomatic Correspondence}, 374–75, 375. Halsey was eventually recalled as consul, and John Quincy Adams as Secretary of State refused to yield to pressures of senators from Rhode Island, Halsey’s home state, to reinstate him. Adams commented in his diary that Halsey “had made himself as obnoxious to the Government at Buenos Ayres, as to his own.” 6 May 1820 entry, The John Quincy Adams Digital Diary, published in the Primary Source Cooperative at the Massachusetts Historical Society: At: \url{https://www.primarysourcecoop.org/publications/jqa/document/jqadiaries-v31-1820-05-p324--entry6?navmode=searchresults&doci=2&ss=obnoxious|Ayres.}\par}
{\scriptsize\color{gray} [p. 10]}\ As a strategy for widening the war, the distribution of blank commissions was very effective. Between 1817 and 1820, with Artigas isolated between Luso-Brazilian invaders and the quietly hostile Buenos Aires government, commissions signed by him were circulating in Baltimore, Savannah, Norfolk, Charleston, and Providence, and cruises against Iberian shipping across the Atlantic and Caribbean were intensifying. When the U.S. government acted to restrict the entry of South American privateers into U.S. ports in 1817, some raiders began offloading booty from Portuguese-flagged vessels at sea. Others took captured ships and cargo to prize courts in the Caribbean, especially at Juan Griego and Angostura off Venezuela, to have the captures declared good prizes. No one knew with certainty whether such tactics would enable South American raiders to be treated as legitimate, licensed privateers in U.S. ports, but there was no shortage of mariners willing to assume the risk.\par
\Needspace{10\baselineskip}\section*{Raids of the \textit{Irresistible}}
{\scriptsize\color{gray} [p. 10]}\ As a raider, John Daniels was not a beginner. He had acquired experience in boarding and seizing ships while serving as a privateer in the War of 1812.\nobreak\textsuperscript{40} When peace arrived in 1814, Daniels returned to commercial trading, sailing for the merchant house of D’Arcy and Didier. Before the war, Baltimore had been in a favored position in West Indies markets. It was the nearest U.S. port for the shipment of wheat from Pennsylvania farms, and it had strong connections to French and Spanish merchant houses operating in the Caribbean. Possibly investing the dowry of his wife, Daniels built a schooner and undertook commercial voyages for D’Arcy and Didier to Haiti and France.\nobreak\textsuperscript{41}\par{\small 40. Daniels was not especially successful as a wartime privateer. On his first outings, he captained two vessels that were seized by British ships and captured only a handful of small vessels. A fourth cruise yielded valuable prizes, but Daniels ran aground in the waters off Cape May, once unintentionally and once intentionally, to avoid capture by British blockading vessels. Hopkins, “For Flag and Profit,” 392–401, 393–94. In contrast, another Baltimore captain who later turned South American privateer, Joseph Almeida, captured 30 vessels during the war and cleared \$46,000 in profits. Head, “Independence on the Quarterdeck,” 7.\par}{\small 41. Daniels had also traded to Haiti for this firm before the War of 1812. Hopkins, “For Flag and Profit,” 394. Daniels’s wife was born as Eugenia Caze in St. Domingue in 1795. She would have arrived in Baltimore sometime after the first surge of refugees from St. Domingue in 1793, when 1,500 refugees, including about 500 Blacks (a mix of free and enslaved people), traveled to the city on fifty-three ships from Cap François ( Hartridge, Walter Charlton, “The Refugees from the Island of St. Domingo in Maryland,” The Maryland Historical Magazine 38, 2 (1943): 103–122 ). The flow of refugees continued, and a community of French speakers formed in the portside neighborhood of Fell’s Point. See Jones, Martha S., “The Case of Jean Baptiste, un Créole de Saint-Domingue: Narrating Slavery, Freedom, and the Haitian Revolution in Baltimore City,” in Ward, Brian, Bone, Martin, and Link, William A., eds., The American South and the Atlantic World (Gainesville: University Press of Florida, 2013), 104–28.\par}
{\scriptsize\color{gray} [p. 11]}\ Hopeful of regaining their competitive advantage in the West Indies trade, Baltimore merchants instead found themselves confronting a crisis of profitability. In peacetime, larger-capacity British ships were moving into markets previously served by the smaller Baltimore schooners.\nobreak\textsuperscript{42} The pinch on profits sent Daniels and others casting around for alternatives, and news of opportunities in the South American republics arrived on cue. In late 1817, Daniels purchased a new brigantine from a Baltimore shipyard and prepared it for an extensive cruise as a privateer.\par{\small 42. Gary Lawson Browne, \textit{Baltimore in the Nation, 1789–1861} (University of North Carolina Press, 1980), chapter 4. British merchants flooding South American ports with goods were also encountering weak demand, as consumer goods, from pianos to mahogany furniture, sat on the docks and in warehouses of South American ports. D. C. M. Platt, \textit{Latin America and the British Trade, 1806–1914} (London: Adam and Charles Black, 1972), 23–64.\par}
{\scriptsize\color{gray} [p. 11]}\ Daniels was attentive to the need to defend his raiding as lawful, and he began to construct a paper trail that might secure future booty. To accept a privateering commission on U.S. soil from a South American polity would have been a violation of U.S. neutrality laws.\nobreak\textsuperscript{43} Hiding their tracks in various ways, other Baltimore captains had already embarked on privateering cruises with commissions from South America, and Daniels probably learned about the opportunity from his contacts and peers.\nobreak\textsuperscript{44} His plan to sail the newly built ship, the \textit{Vacunia}, to the Río de la Plata may have been intended to support the fiction that he obtained his commission outside the United States. Daniels undertook other precautions, too. He stored the ship’s extensive cache of small arms below deck and listed Tenerife as the voyage destination when the vessel set sail in March 1818. The ship’s papers named John Cox as the captain, and Daniels—no doubt to give himself plausible deniability about planning a privateering voyage in a U.S. port—traveled with additional crew members on a pilot boat to meet the ship at the White Rocks, about 9 miles outside Baltimore harbor.\nobreak\textsuperscript{45}\par{\small 43. The language of neutrality legislation changed in response to the emergence of South American polities of uncertain international standing. The Neutrality Act of 1794 had prohibited U.S. citizens from accepting and exercising commissions to “serve a foreign prince or state in war by land or sea.” In response to pressure from the Portuguese minister Corrêa da Serra in December 1816, President Madison endorsed a change to legislation. The revised act passed in March 1817 broadened this language about aid to foreign entities, prohibiting acts to fit out and arm vessels with the intent to “be employed in the service of any foreign prince or state, \textit{or of any colony, district or people}” to commit or aid in hostilities against “\textit{any colony, district or people}, with whom the United States are at peace” (emphasis added). In April 1818, another version added that the ban applied to warlike actions “by land or sea.” Edward Dumbauld, “Neutrality Laws of the United States,” \textit{The American Journal of International Law} 31, 2 (1937): 258–70; and Arlyck, \textit{The Nation at Sea}, 94–95, 204–5.\par}{\small 44. For example, Joseph Almeida and James Chaytor had already launched privateering cruises from Baltimore (David Head, “Independence on the Quarterdeck.”)\par}{\small 45. Depositions of Samuel Beaver and Lawrence Madison, US v. La Irresistible, Appellate Case Files U.S Circuit Court for the District of Maryland (hereafter USCC Maryland), National Archives and Records Administration at Philadelphia (hereafter NARA Philadelphia). Bernabeu v. Nereyda, Admiralty Dockets and Case Files of the U.S. District Court for the District of Maryland (hereafter AD USDC Maryland), NARA Philadelphia.\par}
{\scriptsize\color{gray} [p. 12]}\ Once at sea, Daniels had the arms brought up and the guns mounted.\nobreak\textsuperscript{46} The \textit{Vacunia} then sprinted, on a three-week, 5,000-mile voyage, to Buenos Aires without calling at another port. When and how Daniels obtained the Artigas commission is a mystery, and it is possible that he already possessed it when he sailed from Baltimore.\nobreak\textsuperscript{47} We know that in Buenos Aires, Daniels went ashore to obtain a commission from the government of the United Provinces. As further steps to support his claim that the subsequent captures were lawful, he recorded a sale and re-purchase of the \textit{Vacunia} and changed its name to the \textit{Maipú}, and became a naturalized citizen of the United Provinces.\nobreak\textsuperscript{48} Daniels then released and rehired his crew, except for a few men who chose to stay ashore or change ships, and he added another twenty men to the ship’s roster. When officials in Buenos Aires became suspicious that Daniels had no intention of bringing prizes back to Buenos Aires, they informed him that they would send marine observers on board. Daniels slipped out of the harbor that night.\nobreak\textsuperscript{49}\par{\small 46. The vessel was said to be armed with “muskets, pistols, cutlasses, and boarding pikes, and also twelve eighteen pound gunnades” (naval artillery), as well as “a quantity of powder, ball and ammunition.” Bernabeu v. Nereyda, AD USDC Maryland, NARA Philadelphia.\par}{\small 47. U.S. court documents contain conflicting claims about the date the commission was issued. One judge stated that the commission from Artigas was dated 14 February 1818—a month before Daniels left Baltimore (Bernabeu v. Nereyda, AD USDC Maryland, NARA Philadelphia). Daniels’s lawyer later claimed that the commission was issued by Artigas at Purificación on 15 June 1818 (J. D. Daniels v. Vasques, USCC Maryland, NARA Philadelphia). In a third case, the Portuguese consul confidently asserted that Daniels had never possessed such a commission (Swift v. J. D. Daniels, AD USDC Maryland, NARA Philadelphia). There is no evidence that Daniels ever presented the commission before U.S. courts.\par}{\small 48. On the fluidity of citizenship in the period, see Pilar González Bernaldo de Quirós, \textit{Argentina hasta la muerte. Políticas de nacionalidad y prácticas de naturalización, siglo XIX-XX} (Buenos Aires: Fondo de Cultura Económica, 2025), chapter 1; Mongey, Vanessa, Rogue Revolutionaries: The Fight for Legitimacy in the Greater Caribbean (Philadelphia: University of Pennsylvania Press, 2020) DOI: 10.9783/9780812297577, 71; and Emily A. Yankowitz, “Who Is a Citizen?: Negotiating American Citizenship Before The Fourteenth Amendment” (PhD diss., Yale University, 2025). On the exclusion of Free Blacks from citizenship, see Jones, Martha, Birthright Citizens: A History of Race and Rights in Antebellum America (Cambridge: Cambridge University Press, 2025).\par}{\small 49. Swift v. J. D. Daniels, AD USDC Maryland, NARA Philadelphia.\par}
{\scriptsize\color{gray} [p. 12]}\ Again at sea, Daniels renamed the ship \textit{La Irresistible}, or the \textit{Irresistible.} He gathered the crew and announced that he possessed two commissions—one from the United Provinces and another from the Oriental Republic—and that he intended to raid under both. A crew member later testified that Daniels showed the crew the commission from the United Provinces, but no one remembered seeing the Artigas commission. Daniels surely knew that prize courts would not approve sailing under two commissions, and perhaps for that reason he decided to abandon the commission from the United Provinces. He sent it back to Buenos Aires by way of a passing ship. In Buenos Aires, officials confiscated the deposit filed for Daniels’s ship by a U.S. merchant and declared Daniels to be a pirate.\nobreak\textsuperscript{50}\par{\small 50. Depositions of Samuel Beaver and Lawrence Madison, US v. La Irresistible, USCC Maryland, NARA Philadelphia. The deposit for the \textit{Maipú} was reportedly paid by William Ford, who would have incurred the loss when Daniels abandoned his Buenos Aires commission. Officials there took the unusual measure of declaring Daniels a pirate because he retained some of the privateering documents (\textit{Gaceta de Buenos Aires} No. 78, 8 julio 1818).\par}
{\scriptsize\color{gray} [p. 13]}\ In the next seven weeks, the \textit{Irresistible} seized twenty-six ships.\nobreak\textsuperscript{51} All but one were sailing under Portuguese colors and plying the busy sea lanes connecting Brazil to Lisbon. In most cases, the crew plundered the cargo and sank the vessels.\nobreak\textsuperscript{52} The richest prizes came from the two most profitable segments of Portuguese Atlantic commerce. One was the \textit{Globo}, a Portuguese-flagged merchant vessel sailing from Bombay to Lisbon with a cargo of cotton, tea, ginger, cinnamon, saltpeter, and other spices. The \textit{Irresistible} seized the cargo, later valued at \$90,000 by the Portuguese consul in Baltimore, and sent the ship, valued at \$30,000, to Charleston, South Carolina, presumably to be sold.\nobreak\textsuperscript{53} A second, even richer capture was of a Portuguese vessel called the \textit{Grão Para} (later referred to in U.S. courts as the \textit{Gran Para}), on a voyage from Rio de Janeiro to Lisbon. The ship was carrying “Fifty Bags of Silver, Ten Bags of Gold and Silver, [and] Five Bags of Gold” valued at about \$300,000.\nobreak\textsuperscript{54} Confident that his privateering commission would protect him from legal challenges, Daniels returned home and placed his share of the captured treasure in the Marine Bank of Baltimore. But he was underestimating how quickly and decisively the political winds in the United States had shifted against him during his months away.\par{\small 51. Some sources list the number as 29 (Gazeta da Lisboa, 27 octubre 1818).\par}{\small 52. The Spanish vessel \textit{Flor de Santa Catalina} was sent with a small crew to the prize court at Juan Griego on the Island of Margarita. Beraza, \textit{Corsarios de Artigas}, 113.\par}{\small 53. Swift v. J. D. Daniels, AD USDC Maryland, NARA Philadelphia.\par}{\small 54. The estimates are from the Portuguese consul’s libel (J. D. Daniels v. Vasques, USCC Maryland, NARA Philadelphia). It is possible that some or all of the treasure was payment for the sale of enslaved Africans carried to Rio de Janeiro. British efforts to prohibit the slave trade by signing bilateral treaties had left a gaping hole in the abolition regime by allowing the trade to continue in Portuguese hands south of the equator. Benton and Ford, \textit{Rage for Order}, chapter 5; see also Richards, Jake Subryan, The Bonds of Freedom: Liberated Africans and the End of the Slave Trade (New Haven: Yale University Press, 2025), chapter 6.\par}
\Needspace{10\baselineskip}\section*{Friendly and Unfriendly Courts}
{\scriptsize\color{gray} [p. 13]}\ U.S. government backing for citizens engaged in raiding for South American polities was fast eroding. Officials were facing intense lobbying by Iberian ambassadors and rapidly escalating litigation by foreign consuls.\nobreak\textsuperscript{55} The government was also immersed in negotiations with the Spanish government over the acquisition of Florida and other terms of what would become the Adams-Onís Treaty, and placating Spain was now a high priority for Secretary of State John Quincy Adams. Initially, Daniels responded to these dangers by quickly preparing the \textit{Irresistible} for another privateering voyage.\par{\small 55. For an excellent analysis of foreign consuls and their role in litigation, see Arlyck, Kevin, “Plaintiffs v. Privateers: Litigation and Foreign Affairs in the Federal Courts, 1816–1822,” Law and History Review 30, 1 (2012): 245–78 DOI: 10.1017/S0738248011000666.\par}
{\scriptsize\color{gray} [p. 13]}\ In his second cruise, Daniels seized several dozen more ships. They included the \textit{Paquete do Maranhão}, a Portuguese slave ship, and a sixteen-gun Spanish warship, the brig \textit{Nereyda.} Daniels recognized that the ship was valuable, and he decided to take it to a prize court to have it declared good prize. He first sought a quick ruling in St. Thomas but was rebuffed by officials there. The \textit{Irresistible} then accompanied the \textit{Nereyda}, with a prize crew made up of Daniels’s men, to Juan Griego on the island of Margarita, where the friendly prize court would recognize his Artigas commission and award the ship to him as lawful captor.\nobreak\textsuperscript{56}\par{\small 56. La Nereyda, 21 U.S. 108 (1823); Beraza, \textit{Corsarios de Artigas}, 115.\par}
{\scriptsize\color{gray} [p. 14]}\ The court at Juan Griego was one of a handful of Caribbean prize courts profiting from the surge in the region’s maritime raiding. Officials in Gran Colombia had created the forum to rule on captures made by privateers sailing for the new government, and the court quickly signaled that its doors were open to captains with commissions from other Latin American states.\nobreak\textsuperscript{57} Following swift action by the court, Daniels conspired with his close associate Henry Childs to arrange what was likely a sham sale of the vessel to a Venezuelan citizen named Antonio Julio Franceschi. The \textit{Nereyda} was rechristened the \textit{Congreso de Venezuela}, and the new owner—with Childs and Daniels acting in the background—secured a commission to sail as a privateer for the revolutionary government of Simón Bolívar.\par{\small 57. Several Caribbean islands under neutral empires—St. Bartholomew (Sweden), St. Eustatius (the Netherlands), and St. Thomas (Denmark)—housed prize courts and welcomed, though inconsistently, South American privateers. Mulich, Jeppe, “Prize court politics and regional ordering in the Caribbean,” Past \& Present, 265, 17 (2024): 193–213 10.1093/pastj/gtae027. On privateering from Venezuela, see Morales, Edgardo Pérez, No Limits to Their Sway: Cartagena’s Privateers and the Masterless Caribbean in the Age of Revolutions (University of Vanderbilt Press, 2018) DOI: 10.2307/j.ctv16757xb; and Mongey, \textit{Rogue Revolutionaries.}\par}
{\scriptsize\color{gray} [p. 14]}\ A mutiny then disrupted Daniels’s plans. The \textit{Irresistible} was anchored in the harbor at Juan Griego when another Baltimore privateering vessel, the \textit{Creola}, sailing with a commission from Buenos Aires, entered the harbor with captures of its own. The crews mingled, and with men of the \textit{Creola} apparently in the lead, they decided to sail for home. The mutineers permitted nineteen men from the \textit{Irresistible}, including Henry Childs, to go ashore. Some of those who remained would later claim they did so under duress. After rifling through the papers on the \textit{Irresistible} in search of a commission and finding none, the merged crew decided to make captures without legal cover. On the way to the Chesapeake, they boarded and ransacked about thirty vessels. Like Daniels, the men discovered on their return that they were entering a political atmosphere that had turned sharply hostile to would-be privateers. Unlike Daniels, they had no claim of a commission and would not be able to argue in court that their captures were lawful acts of war. Officials rounded up crew members and put separate groups on trial for piracy in Baltimore and Richmond.\nobreak\textsuperscript{58}\par{\small 58. \textit{Niles Weekly Register}, 7 August 1819, 390–95.\par}
{\scriptsize\color{gray} [p. 14]}\ Daniels managed to take advantage even of this calamity. After the mutiny at Juan Griego, he commandeered the \textit{Congreso de Venezuela} (the renamed \textit{Nereyda}) in pursuit of the \textit{Irresistible} and the mutineers. Arriving in Baltimore, he learned that the \textit{Irresistible} was docked in Nottingham and that many of the mutineers were under arrest in Richmond. Daniels traveled to retrieve the ship, and port officials readily gave it to him. Once again, he found himself in possession of the \textit{Irresistible}, and now he also had command of the rechristened \textit{Nereyda.} Yet Daniels’s position had never been less secure, and he found himself badly in need of a lawyer.\par
{\scriptsize\color{gray} [p. 14]}\ He was facing three claims, called libels, to recover captured ships and cargo. The Portuguese consul to the United States, Joaquim José Vasques, had sued for the return of the treasure of the \textit{Gran Para.} The Portuguese consul in Maryland, William Swift, presented a libel on behalf of the Portuguese crown to recover the value of goods seized from the \textit{Globo} and twenty-five other vessels on the first cruise of the \textit{Irresistible.} And the Spanish consul in Baltimore, Juan Bautista Bernabeu, demanded that Daniels return the \textit{Nereyda} to the Spanish government. The cases appeared before Theodorick Bland, a Maryland judge known for his support of South American revolutions and at one time suspected of involvement in privateering.\nobreak\textsuperscript{59} The cases would spotlight the shifting contours of the emerging Atlantic system by generating debates over the definition of states, war, and piracy.\par{\small 59. Bland was rumored to have been a member of the American Concern, a shadowy group that formed in Baltimore in 1816 to organize South American privateering ventures. Daniels’s employers, the merchants John D’Arcy and Henry Didier, were also members (Charles C. Griffin, “Privateering from Baltimore during the Spanish American Wars of Independence,” \textit{Maryland Historical Magazine}, XXXV, 1940, 1–25, 3; Harold A. Bierck, “Spoils, Soils, and Skinner,” \textit{Maryland Historical Magazine} XLIX, 1954, 21–40, 28). Bland had served on the president’s Special Commission to Chile and Buenos Aires, dispatched in 1817 to report on the political situation in South America. Bland drew criticism for traveling on his own to Santiago, Chile, in pursuit of a promise from the new Chilean government to assume the debt of an ambitious Chilean revolutionary, José Miguel Carrera, to Bland’s son-in-law, John Stuart Skinner, the Baltimore postmaster. In his report, Bland criticized the government of the United Provinces and expressed some sympathy with the breakaway provinces of the Banda Oriental and Entre Ríos (Stewart Watt, “The South American Commission, 1817–1818,” \textit{Hispanic American Historical Review} 9, 1 [February 1929]: 31–59).\par}
\Needspace{10\baselineskip}\section*{What is a State?}
{\scriptsize\color{gray} [p. 15]}\ The dispute over the capture of the \textit{Gran Para} was on one level straightforward. The case ultimately turned on the question of whether Daniels had violated U.S. neutrality laws by fitting out the ship for war in Baltimore.\nobreak\textsuperscript{60} As judge in the district court, Bland declared himself unconvinced by Daniels’s claim that the ship had been on a commercial voyage before obtaining a privateering commission in the Río de la Plata. The ship had been prepared for war in Baltimore, Bland found, ruling that Daniels had to return the rich booty from the \textit{Gran Para.} His judgment was affirmed on appeal in the circuit court of Maryland. The case next reached the U.S. Supreme Court, where Chief Justice John Marshall pronounced himself certain, too, that “there was no commercial voyage.” Even if the crew had not been formally contracted to join a raiding expedition, Marshall concluded, “Every thing they saw, every thing that was done, spoke a language too plain to be misunderstood.” Daniels had organized—in a port in the United States—a piratical voyage, and nothing that occurred in the Río de la Plata had removed that taint.\nobreak\textsuperscript{61}\par{\small 60. The Gran Para, 20 U.S. 471 (1822).\par}{\small 61. Ibid., 488.\par}
{\scriptsize\color{gray} [p. 15]}\ The \textit{Gran Para} case also prompted wide-ranging debates about the sovereignty of South American governments. Appearing for the Portuguese crown, David Hoffman argued that Daniels could not have been acting as a privateer because “the Banda Oriental is not a state or nation invested with the attributes of sovereignty.” Nor had the U.S. government “recognized Artigas as engaged in a civil war with Spain, or in a war of any kind with Portugal.” Going still further, Hoffman proposed that recognizing Artigas as a head of state would disrupt the entire region: “If Artigas, and his wretched and savage followers be recognized as qualified to wage \textit{war}, then may every township, district, city, village, hamlet, or individual, claim the same high prerogatives.” The Banda Oriental also could not be a state, Hoffman argued, because it did not look or act the part. A state had to possess a “fixed \textit{domain}, a \textit{national treasury}, a \textit{national force}, a \textit{code of laws},” and to make war by sea it needed to have seaports. Hoffman asserted that even if Artigas governed as if he were a head of state, he would remain the leader of only a “horde of men.”\nobreak\textsuperscript{62}\par{\small 62. Ibid., 478. Emphasis in original.\par}
{\scriptsize\color{gray} [p. 16]}\ The possibility that the U.S. government might have already indicated that the Oriental Republic was sovereign lurked in the background. Hoffman acknowledged that in Monroe’s second address to Congress on 17 November 1818, the president had stated, “It appears … that the government of Buenos Ayres declared itself independent in July, 1816 … [and] that the Banda Oriental, Entre Rios, and Paraguay, with the city of Santa Fe, all of which \textit{are also independent}, are unconnected with the present government of Buenos Ayres.”\nobreak\textsuperscript{63} Hoffman could not ignore Monroe’s statement. But he suggested that the president had recognized only some aspects of sovereignty in the Banda Oriental and not “the high privileges of sovereignty.”\nobreak\textsuperscript{64}\par{\small 63. Ibid., 479. Emphasis in original. Note Monroe’s insertion of the qualifier “It appears.”\par}{\small 64. Ibid., 480.\par}
{\scriptsize\color{gray} [p. 16]}\ The crux of this argument was that former colonies possessed only a limited capacity to make war. Hoffman reasoned that a colony recognized by other governments as “engaged in a civil war for the assertion of its independence” might oppose only “the parent country and its allies.” It could not “wage a distinct, and independent war with other nations.” Under this logic, the Banda Oriental was something other than a state among states, having been awarded only “limited and partial” recognition by the United States. The courts could not remedy this situation because they lacked the authority “to declare that a people \textit{do not constitute} a state … [or] that they \textit{are} a state.”\nobreak\textsuperscript{65}\par{\small 65. Ibid., 482. Emphasis in original. Only the executive could make this determination, as the court decided in United States v. Palmer, 16 U.S. 610 (1818).\par}
{\scriptsize\color{gray} [p. 16]}\ In reply, Daniels’s attorney, William Winder, asserted that “the Banda Oriental was a sovereign state \textit{de facto.}” Winder was echoing arguments he had made in the case of the \textit{Santissima Trinidad}, a case that turned on whether a ship’s crew had been augmented illegally in Norfolk, Virginia, before making captures. In that case, Winder had described the act of issuing privateering commissions as “the very highest act of sovereignty” and therefore “peculiarly above ordinary judicial control in foreign countries.” A state at war was by definition a sovereign state; it could not be restricted to making war only to the degree that neutral tribunals would permit.\nobreak\textsuperscript{66}\par{\small 66. The Gran Para, 20 U.S. 471 (1822), 485; The Santissima Trinidad 7 Wheat. 283, 330.\par}
{\scriptsize\color{gray} [p. 16]}\ These arguments spilled beyond the courtroom. As the cases involving Daniels and other captains moved through the system, Iberian diplomats were pressuring U.S. officials to do more to rein in privateering from U.S. ports. Not surprisingly, their arguments echoed those of the lawyers hired and directed by Iberian consuls. In March 1819, the Portuguese ambassador, José Corrêa da Serra, visited John Quincy Adams to demand (in “crippled English,” Adams noted privately) that the government declare all captains and crews sailing on ships with commissions from Artigas to be pirates. The ambassador was fixated on Artigas’s lack of seaports under his command. Corrêa da Serra urged Adams to imagine his reaction if privateers attacked U.S. citizens and their property “under the flag and Commission of Tecumseh.”\nobreak\textsuperscript{67} The complaint added to one months before in which Corrêa da Serra railed that “privateering Commissions from Artigas were mere absurdities” and carried as much weight as commissions “from the Seminole Indians.”\nobreak\textsuperscript{68}\par{\small 67. John Quincy Adams, 19 March 1819 entry, \textit{The John Quincy Adams Digital Diary.} At: \url{https://www.primarysourcecoop.org/publications/jqa/document/jqadiaries-v31-1819-03-p050--entry19.}\par}{\small 68. John Quincy Adams, 17 October 1818 entry, \textit{The John Quincy Adams Digital Diary.} At: \url{https://www.primarysourcecoop.org/publications/jqa/document/jqadiaries-v30-1818-10-p405--entry17.}\par}
{\scriptsize\color{gray} [p. 17]}\ Although in making these analogies Corrêa da Serra was focusing on Artigas’s lack of access to ports, the comparisons to Native American polities were hardly casual. Floating phrases such as “limited and partial recognition” and “\textit{de facto} statehood” and debating the kinds of wars that former colonies could legitimately wage, advocates in privateering disputes were grasping at ways to situate some political communities in an in-between, unsettled status. It would be another fifteen years before Chief Justice Marshall would characterize the Native American polities as “domestic dependent nations.”\nobreak\textsuperscript{69} There was nothing approaching a settled vocabulary to draw on to describe quasi-sovereignty, sovereign indeterminacy, or incompletely recognized states. The ambiguities left South American governments in a bind. By issuing privateering commissions, they were asserting the right of sovereigns to make war but exposing themselves to accusations of sponsoring piracy.\par{\small 69. Cherokee Nation v. Georgia, 30 U.S. 1 (1831).\par}
\Needspace{10\baselineskip}\section*{What is War?}
{\scriptsize\color{gray} [p. 17]}\ The link between state status and rights to make war featured prominently in the second major suit against Daniels, the case of the \textit{Nereyda.} As in the \textit{Gran Para} case, there was strong evidence that Daniels had violated U.S. neutrality laws. In support of the libel by the Spanish consul to recover the \textit{Nereyda}, two seamen from the \textit{Irresistible} confirmed that the ship had left Baltimore carrying arms below deck and fifty or sixty men—more than a typical crew for a commercial voyage.\nobreak\textsuperscript{70} But the case had another element, one that led Judge Theodorick Bland to rule in favor of Daniels. A prize crew from the \textit{Irresistible} had taken the \textit{Nereyda} to the court at Juan Griego, where it was declared a lawful prize. Daniels’s associate Henry Childs gave detailed testimony about the court proceedings and the sale of the \textit{Nereyda} to the Venezuelan citizen Franceschi.\nobreak\textsuperscript{71}\par{\small 70. The seamen also testified that another 20 men—but no additional arms—were taken on board in Buenos Aires—further evidence that the ship was already outfitted for war (Bernabeu v. Nereyda, USCC Maryland, NARA Philadelphia).\par}{\small 71. Bernabeu v. Nereyda, USCC Maryland, NARA Philadelphia.\par}
{\scriptsize\color{gray} [p. 17]}\ Bland backed Daniels. He began by affirming that the \textit{Irresistible} had been outfitted for war in Baltimore in violation of the neutrality act. But, Bland continued, the case before him was not about Daniels’s criminal responsibility, it was about the “seizure by a lawfully commissioned cruiser of a belligerent.” Bland acknowledged widespread uncertainty about the status of nations “born at once” in the southern Atlantic and “the station they were to be allowed to assume” according to “the older nations of the earth.” He regarded these questions as “generally rather political than legal.” Civil war was still a kind of war, Bland observed. If a neutral nation declared a conflict to be a civil war, “its courts cannot consider as criminal those acts of hostility, which war authorizes.”\nobreak\textsuperscript{72}\par{\small 72. Ibid.\par}
{\scriptsize\color{gray} [p. 18]}\ Like the lawyers on both sides of the \textit{Gran Para} case, Bland parsed Monroe’s presidential messages about South American states. He noted that on 2 December 1817, Monroe had declared that the former colonies of Spain were not engaged in an “ordinary insurrection or rebellion” but in a “civil war between parties nearly equal.” Bland quoted Monroe’s second address to Congress in November 1818, in which he named the Banda Oriental as “independent,” and he cited the president’s message of 17 December 1819, which affirmed that U.S. ports were “equally open to both parties.” Monroe had commented that Buenos Aires was steadily maintaining its claim to independence, alongside “the Provinces north of the La Plata bordering on it.” Bland concluded that the president had clearly “designated that country of which José Artigas is the Chief, as one capable of having a maritime force, and as conveniently situated for such a force.” The executive had also explicitly recognized a “civil war between that country and Spain.”\nobreak\textsuperscript{73}\par{\small 73. Ibid.\par}
{\scriptsize\color{gray} [p. 18]}\ Bland now turned to the question of the relationship between the Banda Oriental and Venezuela. President Monroe, he asserted, had identified only “two parties in the civil war, ‘Spain and the Colonies.’” With a confidence that would irk Justice Story when the case arrived on appeal at the Supreme Court, Bland stated without qualification that “the Executive has distinctly recognized those two governments of the Banda Oriental, and of Venezuela as cobelligerents, associates, companions, and allies in the war against Spain.” According to the law of nations, Bland continued, captures by an authorized belligerent’s vessel, even one “illegally fitted out” in the ports of a neutral, were valid. He concluded, “There can be no doubt that the commission of the Irresistible from the government of the Oriental Republic was legal.” On appeal, the circuit court of Maryland affirmed Bland’s ruling that the \textit{Nereyda} did not have to be restored to the Spanish government.\nobreak\textsuperscript{74}\par{\small 74. Ibid. Emphasis in the original.\par}
{\scriptsize\color{gray} [p. 18]}\ When the case reached the Supreme Court, lawyers for the Spanish side tried a new tactic, heaping doubt on the jurisdiction and proceedings of the prize court at Juan Griego and asserting an expansive authority of U.S. courts to evaluate foreign prize proceedings.\nobreak\textsuperscript{75} They began by asserting that there was “no evidence whatever of an alliance between Venezuela and the Banda Oriental.” Even though Monroe had acknowledged that the Banda Oriental and Venezuela were “at war with Spain,” it did not necessarily follow that they were allies—or even engaged in the same war. Nor was Monroe’s statement “sufficient to authorize the Court to allow to Artigas all the rights of war.”\nobreak\textsuperscript{76} Going further, attorneys for Spain proposed that even if Venezuela was recognized as fully independent, its law could not be held to contradict that of the United States since “all sovereignty is, in its nature, as a general rule, local.” U.S. courts’ jurisdiction derived from the principle that a foreign prize court’s proceedings were “always open to inquiry” and its acts could always be nullified by another power.\nobreak\textsuperscript{77}\par{\small 75. La Nereyda, 21 U.S. 108 (1823). The pivot away from questioning the meaning of civil war perhaps reflected the Supreme Court’s decision in \textit{U.S. v. Palmer}, the case of a privateer with a Buenos Aires commission (United States v. Palmer, 16 U.S. 610 [1818]). Marshall wrote that if the United States “remains neutral, and recognizes the existence of a civil war, its courts cannot consider as criminal those acts of hostility which war authorizes” (Rubin, \textit{The Law of Piracy}, 158). See also Arlyck, \textit{The Nation at War}, 217–24.\par}{\small 76. Note that Monroe had named the Banda Oriental, not the Oriental Republic, the confederation listed on Artigas’s commissions.\par}{\small 77. Bernabeu v. Nereyda, USCD Maryland, NARA Philadelphia.\par}
{\scriptsize\color{gray} [p. 19]}\ Justice Story sidestepped these issues in ordering the return of the \textit{Nereyda} to Spain. His ruling focused on the fact that Daniels had not produced documents from the court at Juan Griego. The failure to respond to the court’s order to do so, Story asserted, created suspicion “as to the whole transaction.” Story also credited evidence that the sale to Franceschi had been a ruse, citing Franceschi’s absence in U.S. proceedings and Daniels’s exclusive management of the ship after its arrival in the United States. Implicitly, Story’s ruling affirmed the court’s authority to question foreign prize court proceedings.\nobreak\textsuperscript{78}\par{\small 78. Ibid.\par}
{\scriptsize\color{gray} [p. 19]}\ Like arguments about the capture of the \textit{Gran Para}, the positions taken in the Nereyda case reveal the political stakes of unresolved questions in the campaign against South American privateering. Advocates interpreted statements by the U.S. president about the sovereignty of South American polities in ways that laid bare the ambiguities of their status. Lawyers representing foreign claimants referred to “partial recognition,” ascribed limited sovereignty to former colonies, and warned of the harm to neutrals of unbridled raiding. In the \textit{Nereyda} case, arguments centered on whether South American polities were at war or in a civil war ignited by rebellion. The debates planted doubts not just about the legitimacy of Daniels’s sponsor but also about the relationships among insurgent governments and the composition of the Atlantic order.\par
\Needspace{10\baselineskip}\section*{What is a Pirate?}
{\scriptsize\color{gray} [p. 19]}\ Daniels’s legal jeopardy did not end with these cases. A third suit was brought—and won—for the recovery of the booty seized from the \textit{Globo} and twenty-five other Portuguese ships.\nobreak\textsuperscript{79} Daniels also faced criminal charges. One indictment was sparked by the two leaders of the \textit{Irresistible} mutiny, who accused Daniels of the murder of a British captain during the ship’s second cruise, when shots were fired on a British merchant vessel after it had lowered its flag. This case, too, came before Judge Bland, who cited the fog of war and acquitted Daniels of intentionally causing the captain’s death. The second set of charges formed part of a response in Baltimore to pressures from Washington officials to curtail South American privateering. Daniels was tried for violating the neutrality act. Theodorick Bland again acquitted Daniels, reasoning in part that the captain had acted outside the period covered by the latest revision of the neutrality act. To a frustrated John Quincy Adams, the verdict was based on “a legal quibble,” an example of a cramped “Law-Logic—an artificial system of reasoning, exclusively used in Courts of Justice, but good for nothing any where else.”\nobreak\textsuperscript{80}\par{\small 79. The ruling against Daniels was confirmed by the circuit court and was not appealed to the Supreme Court (Swift v. J. D. Daniels, AD USDC Maryland, NARA Philadelphia.)\par}{\small 80. John Quincy Adams, 26 May 1819 entry, \textit{The John Quincy Adams Digital Diary.} At: \url{https://www.primarysourcecoop.org/publications/jqa/document/jqadiaries-v31-1819-05-p102--entry26.}\par}
{\scriptsize\color{gray} [p. 19]}\ If Daniels escaped punishment, some crew members of the \textit{Irresistible} were not so lucky. Seven mutineers were tried in Baltimore for piracy, and nineteen others were tried in Richmond in July 1819.\nobreak\textsuperscript{81} \textit{United States v. Smith} (named for Thomas Smith, one of the ringleaders) was heard in the circuit court of Virginia, with Chief Justice Marshall serving as one of two justices on circuit, and presented at the U.S. Supreme Court before Justice Joseph Story and Justice Henry Brockholst Livingston. Two influential lawyers argued the case, Attorney General William Wirt for the government and Daniel Webster for the accused.\nobreak\textsuperscript{82}\par{\small 81. There were nine trials conducted for 19 men in Virginia. Samuels, Joel H., “The Full Story of United States v. Smith, America’s Most Important Piracy Case,” Penn State Journal of Law \& International Affairs 1, 2 (2012): 320–62, 348.\par}{\small 82. United States v. Smith, 18 U.S. 153, 183 (1820).\par}
{\scriptsize\color{gray} [p. 20]}\ The piracy case unfolded against the background of efforts to strengthen laws to contain raiding. Faced with the surge in privateering involving U.S. citizens, Congress passed legislation in March 1819—only a month before the return of Daniels and the crew of the \textit{Irresistible} to U.S. ports—to widen the definition of piracy.\nobreak\textsuperscript{83} The act established that anyone on the high seas who committed “the crime of piracy, \textit{as defined by the law of nations}” could be “punished with death” if convicted.\nobreak\textsuperscript{84} It did not specify how the courts would determine the definition of piracy under the law of nations.\par{\small 83. Part of the difficulty stemmed from recent attempts to strengthen laws against piracy, precisely in response to the creative strategies of North American mariners engaged in raiding. Besides the authorization in the Constitution of Congress to make laws to regulate piracy, the baseline for policy was congressional legislation in 1790, the Act for the Punishment of Certain Crimes Against the United States (Act of Apr. 30, 1790, ch. 9, 8, 1 stat. 112 [1790]). The Act established that murder, robbery, and other crimes on the high seas would be punishable by death if the same crimes committed on land were defined as capital crimes under the laws of the United States. In 1812, a Massachusetts case raised the question of whether piracy had to include violence (United States v. Tully, 28 F. Cas. 226 [C.C. Mass 1812]; Alfred P., Rubin, \textit{The Law of Piracy}, International Law Studies 63 [Newport: Naval War College Press, 1988], 128–137). In 1818, the U.S. Supreme Court interpreted the 1790 act to mean that a robbery committed on the high seas “belonging exclusively to subjects of a foreign state” could not be considered piracy and punishable by death if the same crime could not carry a death sentence in U.S. courts (quoted in Samuels, “The Full Story of United States v. Smith,” 334).\par}{\small 84. Act to Protect the Commerce of the United States, and Punish the Crime of Piracy, Act of Apr. 30, 1790, chapter 9, 8, 1 stat. 112 (1790). Emphasis added.\par}
{\scriptsize\color{gray} [p. 20]}\ The criminal case against the mutinous crew of the \textit{Irresistible} tested interpretations of the piracy statute.\nobreak\textsuperscript{85} The jury had no trouble reaching the conclusion that the mutineers had seized the \textit{Irresistible} at Juan Griego and “proceeded to sea on a cruize, without any documents of commission whatever.” They also found that the crew had captured and plundered a Spanish vessel, as charged. But jurors could not decide whether such an act amounted to piracy under the law. With Marshall’s encouragement, they issued a special verdict: “If the plunder and robbery aforesaid be piracy under the act of the Congress of the United States … then we find the said prisoner guilty; if the plunder and robbery, above stated, be not piracy under the said act of Congress then we find him, not guilty.”\nobreak\textsuperscript{86} The question that baffled the jury and divided the judges was whether the defendants had committed piracy under the law of nations and could therefore be punished under the act of Congress.\par{\small 85. The 1819 statute was revised and strengthened in 1820. See note 96 below.\par}{\small 86. United States v. Smith, 18 U.S. 153, 183 (1820), 153–154.\par}
{\scriptsize\color{gray} [p. 20]}\ Attorney General William Wirt declared that there was no mystery. All “writers on the public law,” he stated, agreed that piracy was “depredation on the seas, without the authority of a commission, or beyond its authority.” The definition was “certain, consistent, and unanimous.” Wirt reasoned that since all nations could try and punish pirates as \textit{hostes humani generis}, it followed that piracy’s definition was securely rooted in the law of nations.\nobreak\textsuperscript{87} Daniel Webster countered that there was no fixed definition of piracy among writers on public law. The purpose of the act was in fact to make Congress responsible for resolving this question.\par{\small 87. Ibid., 155.\par}
{\scriptsize\color{gray} [p. 21]}\ In the Supreme Court, Justice Story agreed with Wirt that there was broad consensus among writers on the law of nations that piracy was robbery at sea outside a state of war. It was “the general practice of all nations” to punish such actions, a pattern that amounted to “conclusive proof” that the definition need not rest on “particular provisions of the municipal code, but upon the law of nations.”\nobreak\textsuperscript{88} In dissent, Livingston echoed concerns raised in the circuit court by Marshall that piracy was punishable only as it was defined under domestic law. The burden was on Congress, Livingston argued, “to select from sources … and then to declare, and with reasonable precision to define, what act or acts should constitute this crime.”\nobreak\textsuperscript{89} Piracy was not like murder or robbery—crimes in every state of the union for which “no one is allowed to allege his ignorance.”\nobreak\textsuperscript{90} Citizens were not so familiar with the law of nations. Further, Livingston observed, it would be unjust to declare “that all offences against the law of nations, without defining any one of them, should be punished \textit{with death.}”\nobreak\textsuperscript{91}\par{\small 88. Ibid., 162.\par}{\small 89. Ibid., 180.\par}{\small 90. Ibid., 182.\par}{\small 91. Ibid., 183. Emphasis added.\par}
{\scriptsize\color{gray} [p. 21]}\ Despite the court’s affirmation of their guilt, the men did not hang. Responding to pressures from port cities for leniency, the government pardoned two men and issued a reprieve for the others.\nobreak\textsuperscript{92} The calls for mercy highlighted the fact that the mutineers’ crimes had not been very violent. The defense presented testimony that the piratical crew had treated its victims gently and profited only slightly from their predations. One crew member testified he was an unwilling participant who had not been able to get to shore after the \textit{Irresistible} was seized in Juan Griego. The press, the jury, and even some judges found it difficult to equate an act of small violence by members of settled port communities with the imagination of unfettered violence by enemies of all.\nobreak\textsuperscript{93}\par{\small 92. Calls for mercy came from the press, clergy, judges, and residents. In April, federal officials received a petition signed by 4,000 residents of Baltimore calling for a pardon of two convicted mutineers (John Quincy Adams, 11 April 1820 entry, \textit{The John Quincy Adams Digital Diary,} at: \url{https://www.primarysourcecoop.org/publications/jqa/document/jqadiaries-v31-1820-04-p303--entry11).} The pardons and reprieves in March were part of a broader response. At a meeting on 13 March 1820, cabinet members discussed whether to issue pardons for forty-five men convicted of piracy and under sentence of death in Baltimore, Richmond, Charleston, Savannah, and New-Orleans. All had sailed with ships carrying South American commissions. The conclusion was to execute some from each group and issue reprieves for the rest. John Quincy Adams, 13 March 1820 entry, \textit{The John Quincy Adams Digital Diary}. At: \url{https://www.primarysourcecoop.org/publications/jqa/document/jqadiaries-v31-1820-03-p274--entry13.} Head estimates that between 1816 and 1820, 129 crew members on South American privateers were charged with piracy, and only seven were executed in Boston, Baltimore, Charleston, and Savannah. Head, David, “Baltimore Seafarers, Privateering, and the South American Revolutions, 1816–1820,” Maryland Historical Magazine 105 (2008): 269–93, 282.\par}{\small 93. Samuels, “The Full Story of United States v. Smith,” 345–46, 355–59; \textit{Niles Weekly Register}, 17 June 1817, 275–76.\par}
{\scriptsize\color{gray} [p. 22]}\ It is doubtful that the case against the \textit{Creola} and \textit{Irresistible} mutineers was, as one scholar claims, the “most important piracy case” in U.S. history and had “a profound effect on the development of international law in the United States.”\nobreak\textsuperscript{94} The case did simplify U.S. courts’ approach to piracy cases by affirming that the law of nations could be incorporated, without elaboration, in a U.S. criminal statute.\nobreak\textsuperscript{95} But the ruling merely shifted from Congress to the courts the burden of developing a body of law on which to rely in characterizing the law of nations.\nobreak\textsuperscript{96} In that task, Story’s opinion was of questionable help. His seventeen-page footnote on definitions of piracy showed that legal writers recognized piracy as a crime under the law of nations, but it did not describe other aspects of “the general practice of all nations” nor specify the connections to piracy under the common law.\nobreak\textsuperscript{97} As Rubin notes, the approach left open the possibility that the label of “piracy” might be used “to condemn whatever crimes the municipal law system of the United States wanted to attach to the label.”\nobreak\textsuperscript{98}\par{\small 94. Samuels, “The Full Story of United States v. Smith,” 320. On other important piracy cases, see Arlyck, \textit{The Nation at Sea}, 202–227; and LaCroix, \textit{The Antebellum Constitution}, 88–104.\par}{\small 95. The case had immediate effect on prosecutions by appearing to break “a logjam” of piracy cases. Rubin, \textit{The Law of Piracy}, 146.\par}{\small 96. I thank David Golove and Kevin Arlyck for helpful comments on this case.\par}{\small 97. United States v. Smith, 18 U.S. 153, 183 (1820), 163–180; and see Rubin, \textit{The Law of Piracy}, 145–146.\par}{\small 98. Rubin, \textit{The Law of Piracy}, 147. The court in fact was called on to test the meaning of piracy just a few years later, in another case involving a Baltimore ship with a commission from Artigas (The Antelope, 23 U.S. 66 1825). The ruling affirmed a lower court judgment that enslaved Africans seized from a Spanish ship and brought to the United States should be returned to the Spanish claimants. Echoing a British ruling by Lord Stowell in the \textit{Le Louis} case in 1816, Marshall found that the slave trade was not illegal under the law of nations since some nations continued to regard the trade as lawful. It did not matter that the 1820 revision of the U.S. piracy statute explicitly defined slave trading as piracy. See Arlyck, \textit{The Nation at Sea}, 256–65; John Thomas Noonan, \textit{The Antelope: The Ordeal of the Recaptured Africans in the Administrations of James Monroe and John Quincy Adams} (University of California Press, 1977); and Lauren Benton, “Abolition and Imperial Law, 1780–1820,” \textit{Journal of Imperial and Commonwealth History} 39, 3: 355–74. In an offshoot of the \textit{Antelope} affair, a jury acquitted the captain of piracy because he was carrying a commission from Artigas (Noonan, \textit{The Antelope}, 53).\par}
\Needspace{10\baselineskip}\section*{States without Recognition}
{\scriptsize\color{gray} [p. 22]}\ As the cases arising from the raids of the \textit{Irresistible} were moving through U.S. courts, questions about which entities would gain recognition as states were playing out in an intricate politics of recognition on both sides of the Atlantic. Leaders of new and emerging polities angled to position their states in the international order by demonstrating their sovereignty.\nobreak\textsuperscript{99} Established empire-states assembled ad hoc policies of recognition that officials hoped would not jeopardize imperial interests or destabilize existing alliances.\par{\small 99. Gaffield, Julia, Haitian Connections in the Atlantic World: Recognition after Revolution (University of North Carolina Press, 2015) DOI: 10.5149/northcarolina/9781469625621.001.0001; Gould, \textit{Among the Powers of the Earth.}\par}
{\scriptsize\color{gray} [p. 22]}\ The transition from a collection of political experiments to a regional system of states gained momentum when, in 1822, the U.S. government recognized five new nations: the United Provinces, Chile, Gran Colombia, Peru, and Mexico.\nobreak\textsuperscript{100} Britain soon followed suit, taking steps toward partial recognition by dispatching consuls to South America in 1823 and beginning negotiations on commercial treaties, then recognizing the United Provinces, Colombia, Mexico, and Brazil in 1825.\nobreak\textsuperscript{101} Behind these events lay a complex imperial project to guide regional political change through a combination of diplomacy and legal maneuvering, delivered through prize litigation and port politics.\par{\small 100. Portuguese officials moved first to recognize Latin American states in 1821. Rojas, Daniel, “The Recognition of Latin-American Independences: A Major Transformation in the History of the Law of Nations,” in Pezé, Warren and Rojas, Daniel E., eds. International Recognition. A Historical and Political Perspective (Tübingen: Mohr Siebeck, 2022).\par}{\small 101. , Matthew McCarthy, , Privateering, Piracy and British Policy in Spanish America, 1810–1830 (Woodbridge: Boydell Press, 2013).\par}
{\scriptsize\color{gray} [p. 23]}\ Consider the Spanish reaction to U.S. privateering. Spanish consuls in U.S. ports devotedly gathered information about privateering activities and organized legal action. They also flooded the Spanish ambassador’s office with reports of captures of Spanish ships and complaints about the audacity of captains and crews returning to U.S. cities without fear of consequences for rampant raiding. The correspondence of Juan Bautista Bernabeu, the Spanish consul in Baltimore between 1795 and 1821, reveals the scope and energy of this consular activism.\nobreak\textsuperscript{102} Bernabeu’s early reports listed the comings and goings of Spanish vessels in the harbor, recorded disease outbreaks, and passed along snippets of local news. In the final five years of his tenure, Bernabeu’s attention was almost entirely absorbed by privateering—or piracy, as Spanish officials described U.S. mariners’ raiding. He sent detailed, sometimes daily missives to the Spanish ambassador and hectored him about funds needed to pay lawyers to represent the Spanish government in privateering cases and support Spanish mariners arriving in Baltimore who had been serving on crews of captured vessels. He described in minute detail the progress of the numerous legal challenges he was sponsoring for the recovery of Spanish property.\par{\small 102. On other Spanish consular actions, see Arlyck, “Plaintiffs v. Privateers”; and Sean Perrone, “John Stoughton and the “Divina Pastora” Prize Case, 1816–1819," \textit{Journal of the Early Republic} 28, 2 (2008): 215–242.\par}
{\scriptsize\color{gray} [p. 23]}\ Bernabeu reserved special disdain for “the pirate Daniels.” He was incensed that Daniels returned with the \textit{Nereyda} as a prize and was “strolling brazenly through the streets” while the lawyer hired to defend Spanish interests was struggling to get the vessel impounded and Daniels arrested for piracy.\nobreak\textsuperscript{103} Even after the marshal took charge of the \textit{Nereyda}, Bernabeu was outraged that Daniels had been able to refit the ship and add to its armaments to make it ready to cruise again, and he fretted that lax port controls might allow Daniels to slip out of the harbor, as other privateers had done.\nobreak\textsuperscript{104} When Daniels was finally tried for piracy, in May 1819, Bernabeu met with lawyers to strategize about the case and attended all the sessions of the trial. Anticipating acquittal, he complained bitterly to the Spanish ambassador about the biased proceedings and predicted that Daniels would soon be “free and white as a dove” and the Spanish side thoroughly “mocked.”\nobreak\textsuperscript{105} After Daniels was acquitted and as the \textit{Nereyda} suit proceeded on appeal, Bernabeu continued to track the captain’s activities, reporting in 1822 that Daniels had returned to Baltimore to recruit men for raiding in the Caribbean as an agent of Simón Bolívar.\nobreak\textsuperscript{106}\par{\small 103. Bernabeu to Onís, 16 April 1819, Archivo General de Administración (hereafter AGA) 54, 07722. All quotes from Bernabeu’s correspondence are my translation.\par}{\small 104. Bernabeu to Purviance, 23 August 1819, AGA 54, 07722.\par}{\small 105. Bernabeu to Mateo de la Serna, Chargé des Affaires, Philadelphia, 21 May 1819, AGA 54, 07722.\par}{\small 106. Bernabeu to Andraga, 3 July 1822, AGA 54, 07722.\par}
{\scriptsize\color{gray} [p. 24]}\ Daniels’s exploits attracted so much attention partly because they involved the brazen capture of a Spanish warship. In other respects, his actions were unexceptional—as was Bernabeu’s response. At the same time he was tracking Daniels, Bernabeu was attending to dozens of other cases involving privateering voyages from Baltimore, many led by captains carrying commissions from polities in the Río de la Plata. In December 1819, Bernabeu reported that there were eight other ships in the harbor subject to ongoing litigation.\nobreak\textsuperscript{107} On a single day in that month, the district court heard matters involving four captures, including a motion about the \textit{Nereyda} and one involving a ship captured by Joseph Almeida, a captain who had cruised with a commission from Artigas and was even more adept than Daniels at evading attempts to block his continued raiding.\nobreak\textsuperscript{108}\par{\small 107. Bernabeu to Mateo de la Serna, 3 December 1819, AGA 54, 07722.\par}{\small 108. Bernabeu to de la Serna, 27 December 1819, AGA 54, 07722.\par}
{\scriptsize\color{gray} [p. 24]}\ Spanish demands to curb privateering and punish U.S. mariners as pirates formed the backdrop to negotiations over the Adams-Onís Treaty, even after negotiations were concluded in 1819. U.S. officials sometimes pushed disputes into the courts to lessen diplomatic tensions. When Luis de Onís, Spanish Minister in the United States, complained directly to Adams about Daniels’s arrival in Baltimore on a captured Spanish warship and demanded the ship’s return, Adams instructed him to have the Spanish consul file a libel for the return of the ship.\nobreak\textsuperscript{109}\par{\small 109. John Quincy Adams, 24 April 1819 entry, \textit{The John Quincy Adams Digital Diary.} At: \url{https://www.primarysourcecoop.org/publications/jqa/document/jqadiaries-v31-1819-04-p076--entry24.}\par}
{\scriptsize\color{gray} [p. 24]}\ British policy, meanwhile, was also developing in response to both geopolitics and specific disputes about privateering. In September 1821, a Spanish privateer had captured the \textit{Lord Collingwood}, a British brig carrying trade goods from Buenos Aires. The vessel was taken to Puerto Rico and condemned as good prize there. Initially, the British government assumed that the reason for the seizure was that the vessel had run a blockade, but it soon emerged that the prize had been condemned on the grounds that it was a foreign vessel trading in Spanish territory, an act prohibited under the Laws of the Indies. This logic confirmed merchants’ fears that the restoration of Fernando VII to the Spanish throne would bring a revival of Spanish power and the return of a closed imperial system. Profiting from a surge in exports, merchant houses were determined to maintain access to South American ports.\nobreak\textsuperscript{110} British officials threatened naval intervention to resolve tensions over the capture, and the Spanish government relented. Foreign Minister George Canning downplayed the tensions with Spain over the affair, describing the resolution as “a local remedy to a local grievance.”\nobreak\textsuperscript{111}\par{\small 110. British exports to Latin America increased from 2.3\% of British exports in 1806 to 12.6\% in 1826 (Rory M. Miller, \textit{Britain and Latin American in the Nineteenth and Twentieth Centuries} [London: Longmans, 1993], 73). On the inflated expectations about the promise of South American markets after their opening, see Platt, \textit{Latin American and British Trade}, 23–64.\par}{\small 111. McCarthy, \textit{Privateering, Piracy and British Policy in Spanish America}, 101–17, 108.\par}
{\scriptsize\color{gray} [p. 24]}\ The \textit{Lord Collingwood} case strengthened the resolve of British officials to craft a policy of gradual recognition. The goal was to maintain Spain as an ally and to thread a course between mollifying Whig-supported commercial elites and placating counterrevolutionary powers.\nobreak\textsuperscript{112} The result was an ad hoc progression of steps toward partial recognition. Soon after news arrived in May 1822 of United States acts of recognition of Spanish states, the Board of Trade presented a bill, endorsed by the foreign secretary Lord Castlereagh, to admit vessels from new polities in South America to British ports. The Permanent Under-Secretary of State for Foreign Affairs, Joseph Planta, explained that the intent of the act was to “make as little fuss about it as we can” so as not to disturb relations with Spain.\nobreak\textsuperscript{113} In the summer of 1822, Castlereagh outlined steps to strengthening \textit{de facto} recognition through the appointment of British consuls to Spanish America and the negotiation of commercial treaties.\nobreak\textsuperscript{114}\par{\small 112. The European politics of counterrevolution and European ambivalence about South American intervention were on bright display in Paris in 1817, when the Spanish government asked Allied ministers to condemn the Portuguese invasion of the Banda Oriental. See Beatrice de Graaf, \textit{Fighting Terror After Napoleon: How Europe Became Secure After 1815} (Cambridge: Cambridge University Press, 2020), 428–438.\par}{\small 113. Planta to Straford Canning, 11 May 1822, quoted in Sanchez, Mariano, \textit{The Woodbine Parish Report on the Revolutions in South America (1822)} : The Foreign Office and Early British Intelligence on Latin America (Liverpool: Liverpool University Press, 2022), 37 .\par}{\small 114. Webster, C. K., “Castlereagh and the Spanish Colonies II. 1818–1822,” The English Historical Review 30, 120 (1915): 631–45 DOI: 10.1093/ehr/XXX.CXX.631. An agreement in 1823 declared British American ports under Spanish control open to British trade, ended attacks on British vessels, and established a claims commission to address compensation. De Graaf shows that Castelreagh was also rejecting appeals for a coordinated European response (\textit{Fighting Terror After Napoleon}, 434–435).\par}
{\scriptsize\color{gray} [p. 25]}\ There were other signs, too, that the British government was prepared to quietly endorse a regional system of states. Castlereagh at one point proposed that European powers cultivate the formation of constitutional monarchies across South America. The proposal went nowhere: the French would not endorse it, Spanish officials would not brook independence in any form, and events in the Americas were in any case quickly outpacing such schemes. But Castlereagh’s dream of a continent of constitutional monarchies shows that even British officials attuned to counterrevolutionary pressures were anticipating the rise of a new Atlantic system. The question was not whether states would be established in South America but how quickly and in what form—and how they could be made to align with British interests.\par
{\scriptsize\color{gray} [p. 25]}\ Pressuring new states to contain privateering was part of this project of alignment. Much as Spanish diplomats were urging the United States government to ban privateering from U.S. ports, the British government sought to pressure South American officials to temper their sponsorship of privateering. Unlike Spanish diplomats in the United States, the British could not exercise influence through consuls, so their vehicle for delivering pressure on the new governments had to be the British Navy.\nobreak\textsuperscript{115} The Admiralty reported some success through naval channels. In Buenos Aires, the government moved to require privateers to register bonds of higher value before sailing, and the new Chilean government agreed to intervene to protect British vessels trading in Peruvian ports still under Spanish control. Yet the British government refrained from demanding more vigorous measures, instead opting for a system of limited intervention nestled within a formal policy of non-intervention.\nobreak\textsuperscript{116} To label as pirates all raiders licensed by the new governments would have called into question the status of the new entities as states, including their right to make war. That effect would have diminished the very claims of sovereignty that the British government was quietly seeking to reinforce.\nobreak\textsuperscript{117} The new Atlantic system was being constructed as an interpolity regime that would regulate, not eliminate, maritime violence and preserve imperial rights to use force and coercion to calibrate order.\par{\small 115. The policy matched a global strategy of empowering captains to decide on measures short of war and to act as legal agents in places where there was no official or permanent representation by British officials. See Benton, Lauren and Clulow, Adam, “How Not to Possess an Island: Pitcairn and the Legal Circuits of British Empire in the Pacific World,” Annales 7,4 (2025): 605–46.\par}{\small 116. On the system of authorized limited intervention in the Pacific in the same period, see Benton, \textit{They Called It Peace}, chapter 5.\par}{\small 117. See McCarthy’s excellent analysis in \textit{Privateering, Piracy and British Policy}, 75–80, 159.\par}
\Needspace{10\baselineskip}\section*{Beyond Piracy}
{\scriptsize\color{gray} [p. 26]}\ The \textit{Irresistible} raids and the disputes they produced about booty and piracy spanned the crucial years of Atlantic political transformation. Daniels followed other U.S. mariners who began to raid for South American governments of uncertain status in 1816. Just seven years later, as the Supreme Court was ruling in the last of the cases involving the \textit{Irresistible} captures, imperial policies had shifted the political horizon. Empires still dominated, and micro-polities persisted. But some entities, like Artigas’s Banda Oriental, had succumbed to imperial invasion, and others had joined emergent nations or negotiated tenuous arrangements to retain a measure of autonomy within imperial spheres of influence. Even as the precise political contours of the region remained unstable, the new order was reaffirming a narrowing list of political possibilities. Daniels conformed to the new order by taking a commission in the Colombian navy, a position that secured his profits and allowed him a comfortable retirement in Baltimore.\par
{\scriptsize\color{gray} [p. 26]}\ Placing maritime violence and its regulation at the center of regional change alters narratives about international legal order in this period. Historians have shown that a discourse about degrees of civilization became a central focus of efforts to define the membership of new American states in the international legal community.\nobreak\textsuperscript{118} This account is correct as far as it goes. In conflicts over the legitimacy of privateering and the political order behind it, we glimpse a much more varied set of connections between state formation and the new regional order. Privateers, their sponsors, and their challengers were forced to address fundamental questions about the shape of the international order as they engaged in conflicts over the lawfulness of captures at sea.\par{\small 118. Obregón, Liliana, “Between Civilisation and Barbarism: Creole Interventions in International law,” Third World Quarterly 27, 5 (2006): 815–32 DOI: 10.1080/01436590600780136; Coates, Benjamin Allen, Legalist Empire: International Law and American Foreign Relations in the Early Twentieth Century (Oxford University Press, 2016) DOI: 10.1093/acprof:oso/9780190495954.001.0001. On tensions between U.S. hemispheric power and regional multilateralism in international law, see Scarfi, Juan Pablo, The Hidden History of International Law in the Americas: Empire and Legal Networks (Oxford University Press, 2017) DOI: 10.1093/acprof:oso/9780190622343.001.0001.\par}
{\scriptsize\color{gray} [p. 26]}\ It was mandatory for polities to position themselves in relation to the shifting order. The United States was a case in point. Emerging tentatively as a regional power, the country remained weak and buffeted by more powerful empires. It has become commonplace to describe the Monroe Doctrine as a framework for U.S. imperialism, but the muscular version of the doctrine was articulated decades later to justify U.S. territorial expansion and the extension of slavery.\nobreak\textsuperscript{119} In the early decades of the nineteenth century, U.S. leaders endorsed an Atlantic system of states in part as a bulwark \textit{against} the power of empires. To recover this vision is not to deny U.S. ambitions of regional hegemony. It allows us to recognize the dominance of European empires, to identify forces other than revolutionary movements shaping the international Atlantic, and to widen the list of participants engaged in constitutional projects at the intersection of sovereignty and international law.\par{\small 119. Jay Sexton, \textit{The Monroe Doctrine: Empire and Nation in Nineteenth-Century America} (Hill \& Wang, 2012). The Trump administration has flattened this history to assert a “Trump Corollary” to the Monroe Doctrine, interpreting it as an unqualified statement about U.S. hemispheric power (“National Security Strategy of the United States of America,” November 2025, accessed on 11 February 2026, at: \url{https://www.whitehouse.gov/wp-content/uploads/2025/12/2025-National-Security-Strategy.pdf).}\par}
{\scriptsize\color{gray} [p. 27]}\ Conflicts over maritime violence highlight the ideological eclecticism and political improvisation of the age. As with Artigas’s constitutionalism, the rhetoric of liberty was consistent with visions of layered sovereignty and confederation, and revolutionary talk among men like Daniels meshed with their opportunism. We have seen, too, that privateering produced an ocean of unintended consequences. Raiders bolstered bids for recognition by some new states while exposing others to destabilizing external pressures. Privateering practices were equally available to anti-imperial agents, slave traders, and counterrevolutionaries. Indeed, once privateering commissions were in circulation, their itineraries were unpredictable. Commissions issued by Artigas showed up in Atlantic prize courts in the possession of active slave traders even after the Banda Oriental had fallen to Luso-Brazilian forces and Artigas was in exile.\nobreak\textsuperscript{120}\par{\small 120. Besides the \textit{Antelope} (23 U.S. 66 1825), examples of ships in prize proceedings decided after Artigas was in exile include the \textit{Arrogante Barcelones} (20 U.S. 496 1822) and the \textit{Fanny} (22 U.S. 658 1824). Ships with Artigas commissions intercepted with enslaved people range from the \textit{Mouche}, a small vessel captured by a British ship between St. Thomas and Tortola with a few enslaved Africans on board (TNA, CO 318/820) to the \textit{General Artigas}, which arrived at Galveston in 1817 escorting two prizes, one a ship carrying 287 slaves (William C. Davis, \textit{The Pirates Laffite: The Treacherous World of the Corsairs of the Gulf} [New York: Houghton Mifflin, 2006] 320–21, 326).\par}
{\scriptsize\color{gray} [p. 27]}\ More broadly, the \textit{Irresistible} story reminds us of the dangers of taking too literally the rhetoric about pirates as enemies of all. In the early nineteenth century, the discourse gained political force while losing legal precision. As Arlyck observes, nothing about the “conceptual elegance” of piracy law prevented it from being “knotty in application.”\nobreak\textsuperscript{121} But even the conceptual clarity of the term “piracy” was dimmed by the indeterminacies of the Atlantic system. Uncertainties about the character of states and the unequal distribution of the right to make war made the law of piracy an interpretive field, not a firm foundation for marking exclusion from the international order.\par{\small 121. Arlyck, \textit{The Nation at War}, 217.\par}
{\scriptsize\color{gray} [p. 27]}\ These findings matter as pirates are frequently—and often casually—cited as models for other supposedly universal enemies, most recently terrorists (including imagined “narco-terrorists” and “domestic terrorists”).\nobreak\textsuperscript{122} The language of universal enmity implies that any state (and perhaps any private organization) can claim jurisdiction to attack and kill people labeled as enemies of all. Yet, as with piracy and privateering, the conclusion that a group sits outside the international legal order is often based on little more than rhetorical flourish. Some groups labeled as outcasts operate with funding from states and some are endorsed or run by quasi-sovereign entities, while very few groups fit the category of truly independent or roguish actors. As with campaigns against sea raiding, the natural law foundations of the “war on terror” matter much less than the structure and synchrony of state policies.\nobreak\textsuperscript{123} In our search for massive disruptions to the international order, we would do well to shift our attention from criminal rogues to criminal polities.\nobreak\textsuperscript{124}\par{\small 122. On the slipperiness of the discourse, see Li, Darryl, The Universal Enemy: Jihad, Empire, and the Challenge of Solidarity (Redwood City: Stanford University Press, 2019).\par}{\small 123. Scheppele, Kim Lane, “The Deep Dilemma of Evidence in the Global Anti‐Terror Campaign,” in Jackson, Vicki and Fabbrini, Federico, eds., Constitutionalism across Borders in the Struggle Against Terrorism (Northampton: Edward Elgar, 2016).\par}{\small 124. Douglas, Lawrence, The Criminal State: War, Atrocity, and the Dream of International Justice (Princeton: Princeton University Press, 2025); Salaymeh, Lena, “Comparing Islamic and International Laws of War: Orthodoxy, ‘Heresy,’ and Secularization in the Category of Civilians,” The American Journal of Comparative Law 69, 1 (2021): 136–67 DOI: 10.1093/ajcl/avab001.\par}
{\scriptsize\color{gray} [p. 28]}\ With these themes and the story of the \textit{Irresistible} in mind, we can return to the debates about piracy with which we began. Romantic representations of pirates deserve our skepticism not only because they fly in the face of copious evidence about sea raiders’ connections and aspirations. They also encourage a dangerously simplistic, inaccurate view of maritime violence and international ordering. Sea raiders like Daniels were, along with the sea of states in which they navigated, forced to become hyper-political actors. They deployed a mix of well-known strategies, from altering ships’ documents to changing citizenship to forum shopping, while inventing new ways to maneuver through political systems of many moving parts. Sea raiding both pinned states firmly in the crosswinds of imperial power and helped to summon into existence a new interpolity order.\par
\end{document}
